\documentclass[11pt,letterpaper]{article}
\usepackage[T1]{fontenc}
\usepackage[utf8]{inputenc}
\usepackage{lmodern,microtype}
\pdfmapfile{+lm.map}
\pdfmapfile{+cm.map}
\pdfmapfile{+symbols.map}
\usepackage[margin=0.9in,headheight=15pt]{geometry}
\usepackage{amsmath,amssymb,amsthm,mathtools}
\usepackage{booktabs,longtable,array,tabularx,enumitem}
\usepackage{xcolor,fancyhdr,listings,xurl,needspace}
\usepackage[colorlinks=true,linkcolor=blue!35!black,citecolor=blue!35!black,urlcolor=blue!40!black]{hyperref}
\usepackage{bookmark}
\definecolor{ink}{RGB}{24,45,64}
\definecolor{pale}{RGB}{245,247,249}
\setlength{\emergencystretch}{3em}
\setlength{\parskip}{0.3em}
\setlist{itemsep=3pt,topsep=5pt,leftmargin=1.6em}
\newtheorem{theorem}{Theorem}[section]
\newtheorem{lemma}[theorem]{Lemma}
\newtheorem{proposition}[theorem]{Proposition}
\newtheorem{corollary}[theorem]{Corollary}
\theoremstyle{definition}\newtheorem{definition}[theorem]{Definition}
\theoremstyle{remark}\newtheorem{remark}[theorem]{Remark}
\newcommand{\ipra}{\mathsf{iPRA}}
\newcommand{\D}{\mathsf D}
\newcommand{\Q}{\mathsf Q}
\newcommand{\Eq}{\mathsf{Eq}}
\newcommand{\Tr}{\mathsf{Tr}}
\newcommand{\FV}{\operatorname{FV}}
\newcommand{\llbracket}{\mathopen{[\![}}
\newcommand{\rrbracket}{\mathclose{]\!]}}
\newcommand{\file}[1]{\expandafter\nolinkurl\expandafter{\detokenize{#1}}}
\newcommand{\Acal}{\mathcal A}
\newcommand{\Ccal}{\mathcal C}
\lstset{basicstyle=\ttfamily\footnotesize,columns=fullflexible,keepspaces=true,
 breaklines=true,showstringspaces=false,frame=single,rulecolor=\color{black!15},
 backgroundcolor=\color{pale},aboveskip=8pt,belowskip=8pt}
\pagestyle{fancy}\fancyhf{}
\fancyhead[L]{\small A deterministic witness rule for iPRA}
\fancyhead[R]{\small 10 November 2026}
\fancyfoot[C]{\thepage}
\hypersetup{pdftitle={A Deterministic Witness Rule for iPRA: Equivalence, Complete Checkers, and Limits},pdfauthor={Research prepared for Vladimir Reshetnikov}}
\title{\color{ink}\bfseries A Deterministic Witness Rule for iPRA\\[8pt]
\large Equivalence, complete checkers, and the limits of the reduction}
\author{Research prepared for Vladimir Reshetnikov}
\date{10 November 2026}
\begin{document}
\maketitle
\begin{abstract}
The independent choice between case introduction and equation-progress
introduction in the supplied paired-binder calculus is admissibly removed.
Progress is selected exactly when the matrix is a literal equation and
the existential witness is the literal zero term; every other introduction
uses cases. An unused bound variable recovers progress at other witnesses.
The independently certified term $C=I(2,1;v.0)=0$ recovers atomic cases at
zero without presupposing substitution admissibility for the new rule.
The resulting calculus preserves every native sequent through the identity
on formulas, and hence preserves the inherited whole-domain
bi-interpretation and equiconsistency with first-order intuitionistic PRA.

The complete new checker has 555 normalized Python bytes. A separate
implementation of the unchanged Choice-DI calculus has 550 bytes. The
inspected current corpus attains a minimum of 558 bytes, with different
programs attaining 262 tokens and 296 AST nodes. Our two selected programs
do not improve those latter minima. The code reduction comes from a
small dispatch simplification and an inherited invalid-image convention;
the new deterministic rule costs five bytes relative to the 550-byte
control. These contributions are accounted for separately.

This is a verified small reduction, not the requested significant new
arithmetic foundation or a proof of absolute simplicity. The package
contains written proofs, complete source, executable proof normalization,
adversarial controls, the frozen source census, and a compiled report.
The mathematical metatheorems are not Lean or Rocq formalizations.
\end{abstract}
\begin{center}
\fcolorbox{ink}{pale}{\begin{minipage}{0.92\linewidth}
\textbf{Established:} one deterministic introduction pattern, the same
native consequence relation, two smaller complete byte profiles, and a
separate restricted-congruence theorem.\\
\textbf{Not established:} a global minimum, a simultaneous improvement
of all size coordinates, or the requested major foundational breakthrough.
\end{minipage}}
\end{center}
\clearpage
\setcounter{tocdepth}{1}
\hypersetup{bookmarksdepth=2}
\begingroup\small\setlength{\parskip}{0pt}\tableofcontents\endgroup
\clearpage

\section{What changed, and what the comparison actually shows}
\label{sec:overview}

The research target is full first-order intuitionistic primitive recursive
arithmetic, with the symbols of fixed finite primitive-recursive
descriptions, their defining equations, the successor axioms, and
quantifier-free induction in intuitionistic first-order equality logic.
Descriptions are external finite syntax. A number does not range over
arbitrary functions, and no uniformly total interpreter is postulated.
The requirement here includes internally provable interpretation
composites, not merely equiconsistency or agreement on standard numerals.

The selected repository snapshot is
\file{76a2651601d17b34d46753f189ea026a21ce8d29}~\cite{snapshot}.
It contains 17 direct ZIPs in \file{docs/drafts}. Recursive inspection
finds 32 distinct ZIP payloads, 443 distinct Python source byte strings,
and 2,239 source occurrences after archive-content deduplication.
Every source at or below at least one observed boundary was inspected:
36 candidates comprise three complete checkers and 33 supporting utilities.
The full source census and all 443 frozen source byte strings are supplied.
The historical sources are parsed for measurement, not executed by the census.

The current byte minimum is the certified-zero program at 558 bytes.
Its distinct joint profile attains 262 significant tokens and 296 AST nodes.
These are the smallest values \emph{attained in this finite collection}.
They are upper bounds on an unknown global optimum, not mathematical
lower bounds on every possible axiomatization or checker.

\subsection{Three contributions that must not be conflated}

\paragraph{A rule-normalization theorem.}
The new calculus $\mathcal C_{\rm w}$ has one introduction pattern whose
choice is determined by the witness syntax. It retains the same numerical
operations, formula meanings, equality transport, and elimination rules.
The proof in Section~\ref{sec:deterministic-witness} preserves all native
sequents. Its main technical point is to use hypothesis cut, not an
unjustified numerical substitution that could change the syntactic guard.

\paragraph{A small implementation improvement.}
Every selected major triple is truthy, and every generated computation
equation is a nonempty array. A Boolean expression therefore shares their dispatch.
Canonical null markers are replaced by empty lists. This saves six
normalized bytes with no change of native theory. The inherited invalid
image $0j$ saves another two bytes relative to \texttt{None}. Thus the
unchanged Choice-DI theory has a complete 550-byte implementation. Adding
the deterministic witness guard gives the 555-byte implementation of
$\mathcal C_{\rm w}$. The mathematical restriction alone is not credited
with the byte saving.

\paragraph{An independent congruence refinement.}
Appendix~\ref{zcc:section} shows that counter congruence for a different,
counter-dependent recursor needs only seed zero and state-free updates.
The displayed numerical-congruence derivation adds twenty nodes, or five
for a distinct wrapped rule. This result does not implement the measured
scalar checker and is assigned no byte record.

\subsection{Exact whole-program sizes}

\begin{center}\small
\begin{tabular}{lrrrr}
\toprule
Complete program & Bytes & Tokens & AST & Gzip-9\\
\midrule
Supplied certified-zero byte profile & 558 & 264 & 310 & 345\\
Supplied certified-zero joint profile & 598 & 262 & 296 & 350\\
Same byte baseline, invalid image only & 556 & 264 & 310 & 343\\
New unchanged-choice implementation & 550 & 264 & 310 & 342\\
New deterministic-witness implementation & 555 & 266 & 312 & 344\\
New witness implementation, explicit type & 559 & 266 & 312 & 348\\
New witness implementation, explicit list & 563 & 266 & 310 & 349\\
\bottomrule
\end{tabular}
\end{center}

The two smallest new files are \file{code/checker_choice.py} and
\file{code/checker.py}. The last two rows use an ordinary exact-type test
and \texttt{isinstance}, respectively; all four avoid the generic-iterable
malformed-string qualification of the old joint profile. The 550-byte
reduction is $8/558\approx1.43\%$; the 555-byte reduction is
$3/558\approx0.54\%$. No selected program is below all three old
coordinate minima. The readable expansion is separately measured in
\file{evidence/measurements.json}, rather than omitted from accounting or
mistaken for a dependency of the compact file.

Measurements include imports, input loading, scope checking, computation,
logic, and the acceptance verdict. Python and its standard JSON library
are the common runtime boundary used by the comparison programs. Exact
measurements use CPython 3.12.14 and zlib 1.3.2, normalize by
\texttt{ast.unparse(ast.parse(raw)) + LF}, count all \texttt{ast.walk}
nodes, and exclude only comment/layout/encoding tokens from token counts.
No axiom data or validator is hidden outside a measured program.
Gzip is \texttt{gzip.compress(raw, compresslevel=9, mtime=0)}, including
its wrapper; it is an informational figure and not an optimization target.

% Includable manuscript fragment. Requires amsmath, amssymb, amsthm.
% Requires theorem and lemma environments; no project-specific macros.
% Canonical mathematical calculus only. Raw JSON alias soundness is separate.
\section{A deterministic introduction rule selected by the witness}
\label{sec:deterministic-witness}

The independent case/progress choice can be removed from paired
introduction without changing the native consequence relation.  The
replacement selects progress precisely when the matrix is a literal
numerical equation and the supplied witness is the literal zero term.
Every other introduction uses the zero case.  This is a restriction on
finite syntax, not a test of the value denoted by a term.

\subsection{The two precisely specified calculi}

Numerical terms are finite expressions
\[
 t::=x\mid0\mid St\mid I(n,a;v.h),
\]
where only $v$ in $h$ is bound.  Every update is fixed finite syntax with
passive parameters, and all displayed substitutions avoid capture.
For the executable interface we use the conservative closure extension of
the Eq/$D$ language:
\[
 A::=(s=t)\mid D_{x,y}A\mid Q_yA.
\]
The primary Eq/$D$ fragment is recovered by replacing $Q_yB$ with
$D_{u,y}B$ for an unused $u$. The two expressions are interderivable:
paired elimination gives the closure in one direction, and case
introduction with the unused witness $S0$, using its base and QE,
gives the other. The inherited
logical bridge explains this conservative notation.
The ordinary reading of a closure is
$B(0)\to\forall y\,B(y)$; the ordinary reading of $D_{x,y}B$ is
$\exists x\,(B(x,0)\to\forall y\,B(x,y))$.  Neither reading is
changed by the construction below.  Contexts are finite reusable lists;
auxiliary closures may occur as local resources.

The two calculi have the same computation axioms.  The proof works with
the unrestricted iterator equations, and also with the smaller
certified-zero inventory
\begin{align*}
 I(0,a;v.h)&=a
 &&\text{when }a\not\equiv_{\rm syn}0,\\
 I(Sn,a;v.h)&=h[I(n,a;v.h)/v]
 &&\text{when }n,a\not\equiv_{\rm syn}0.
\end{align*}
Here $\equiv_{\rm syn}$ denotes literal syntactic identity.  Computation
axioms may be leaves in every context.  Equivalently, a checker may supply
such an axiom as a local assumption and subsequently use hypothesis
closure; the cut lemma below justifies the equivalence.

The common logical rules are exactly hypothesis use, equality transport,
closure specialization, and paired elimination:
\[
 \frac{}{\Gamma\vdash A}\ (A\in\Gamma),
 \qquad
 \frac{\Gamma\vdash s=t\qquad\Gamma\vdash F(s)}
      {\Gamma\vdash F(t)},
\]
\[
 \frac{\Gamma\vdash Q_yB(y)\qquad\Gamma\vdash B(0)}
      {\Gamma\vdash B(r)},
 \qquad
 \frac{\Gamma\vdash D_{x,y}B(x,y)\qquad
       \Gamma,Q_yB(b,y)\vdash A}
      {\Gamma\vdash A}.
\]
The numerical parameter $b$ in paired elimination is fresh for the
retained context, conclusion, and passive matrix parameters.  The
transport motive $F$ is any fixed native formula template, including
one whose hole lies beneath a binder; binders are first renamed away
from the substituted terms.

Let $\mathcal C$ additionally have the two inherited introductions
\[
 \frac{\Gamma,B(t,0)\vdash B(t,Sz)}
      {\Gamma\vdash D_{x,y}B(x,y)}
 \quad\text{(case)},
 \qquad
 \frac{\Gamma,E(t,z)\vdash E(t,Sz)}
      {\Gamma\vdash D_{x,y}E(x,y)}
 \quad\text{(progress)},
 \tag{DI$_{\rm old}$}
\]
where $E$ must be a literal equation.  In both rules $z$ is fresh for
$\Gamma$, the witness $t$, and every passive matrix parameter.

For the new calculus $\mathcal C_{\rm w}$ define, on fixed syntax,
\[
 \delta_{B,t}(z)=
 \begin{cases}
 z,& B\text{ is a literal equation and }t\equiv_{\rm syn}0,\\
 0,&\text{otherwise}.
 \end{cases}
\]
Its only paired-introduction pattern is
\[
 \frac{\Gamma,B(t,\delta_{B,t}(z))\vdash B(t,Sz)}
      {\Gamma\vdash D_{x,y}B(x,y)}.
 \tag{DI$_{\rm w}$}\label{eq:det-witness-intro}
\]
It has the same freshness conditions.  In particular, a nonliteral
term denoting zero is still a case witness.  Compound matrices always
use cases.  No additional induction, substitution, reflexivity, symmetry,
or matrix-congruence rule is assumed.

\subsection{Structural facts established before the simulation}

\begin{lemma}[Renaming, weakening, and ordinary-assumption cut]
\label{lem:det-witness-cut}
In $\mathcal C_{\rm w}$, injective renaming of numerical variables,
weakening of reusable contexts, and ordinary-assumption cut are
admissible.  The latter has the form
\[
 \frac{\Gamma\vdash A\qquad\Gamma,A\vdash H}
      {\Gamma\vdash H},
\]
where $A$ is an object formula.  The same conclusions hold with additional
retained local closure clauses.
\end{lemma}
\begin{proof}
Variable renaming leaves both tests in
\eqref{eq:det-witness-intro} unchanged: a literal equation remains a
literal equation, and renaming never turns a nonzero syntactic term
into the literal constant $0$, or conversely.  It also preserves the
syntactic conditions on computation axioms.  Structural induction on
a derivation therefore proves renaming admissible.  Injections on the
finitely many variables occurring in a proof suffice; fresh variables
can always be chosen outside another fixed finite derivation.

For weakening, first rename local eigenparameters away from the added
context, then retain that context throughout the original proof.
Every inference is still an instance of the same rule, with the same
witness syntax.  Exchange and contraction are harmless because contexts
are reusable lists and hypothesis closure is membership.

For cut, induct on the derivation of $\Gamma,A\vdash H$.  At a use
of the cut assumption, insert the given derivation of $A$, weakened by
any local assumptions.  Other hypotheses and computation leaves are
unchanged.  Transport and closure specialization commute with this
replacement.  At a paired elimination or introduction, first rename
its eigenparameter away from the finite proof being inserted.  Apply
the induction hypothesis to its premise derivation or derivations and
repeat the original inference.  This operation substitutes derivations
for hypothesis uses; it never substitutes numerical terms into the
witness field.  Consequently the literal-zero test and the
literal-equation test are preserved at every introduction.  This proves
cut without assuming admissibility of arbitrary numerical substitution.
\end{proof}

\begin{lemma}[A closed nonliteral representative of zero]
\label{lem:det-witness-zero}
Both calculi prove $C=0$ and $0=C$, where
\[
 C:=I(S(S0),S0;v.0).
\]
These proofs use only permitted computation axioms and equality
transport, even with the restricted computation inventory above.
\end{lemma}
\begin{proof}
The successor computation with predecessor $S0$ and initializer $S0$
gives $C=0$ directly: both supplied terms are nonliteral zero, and the
constant update returns $0$.

To obtain the reverse equation without presupposing symmetry or
reflexivity, put $K:=I(0,C;v.0)$.  The permitted base computation gives
$K=C$, since its initializer $C$ is not the literal zero term.  Apply
equality transport along $K=C$ to a second copy of $K=C$, with motive
$F(w):=(w=C)$.  The conclusion is $C=C$.  Now transport this conclusion
along $C=0$, with the same motive, obtaining $0=C$.

Thus neither a base computation with literal-zero initializer nor an
introduction rule is used in the construction.  If desired, $0=0$
follows separately by transporting $C=0$ along itself with motive
$w=0$; it is not needed as a hidden premise of the displayed proof.
\end{proof}

\begin{lemma}[Computation repair independent of introduction]
\label{lem:det-witness-computation-repair}
Every unrestricted iterator computation equation is derivable from the
restricted computation inventory by equality transport alone.  No
paired introduction, ordinary-assumption cut, or admissible numerical
substitution is needed for this assertion.
\end{lemma}
\begin{proof}
The theorem $C=0$ was obtained directly in
Lemma~\ref{lem:det-witness-zero}.  A base equation with initializer
literal zero is repaired using the fixed motive
\[
 F(w):=\bigl(I(0,w;v.h)=w\bigr).
\]
Its instance $F(C)$ is a permitted computation axiom.  Transport along
$C=0$ gives $F(0)$, the missing computation.

For a successor computation, replace each forbidden literal-zero
supplied field by the same fresh placeholder $w$.  When only the
initializer is forbidden, use
\[
 F(w):=\bigl(I(Sn,w;v.h)=h[I(n,w;v.h)/v]\bigr).
\]
When only the predecessor is forbidden, use
\[
 F(w):=\bigl(I(Sw,a;v.h)=h[I(w,a;v.h)/v]\bigr).
\]
When both are forbidden, use
\[
 F(w):=\bigl(I(Sw,w;v.h)=h[I(w,w;v.h)/v]\bigr).
\]
In each case $F(C)$ is a permitted successor computation.  Transport
along $C=0$ gives the exact desired equation $F(0)$.  Choose $w$ fresh
for $h$ and all passive parameters, and rename the bound state variable
before forming each motive.  Hence only the displayed argument positions
are changed; unrelated zero occurrences and passive parameters stay fixed.
This is a construction of individual proofs by the primitive transport
rule, not an assumed operation substituting terms into whole proofs.
\end{proof}

\subsection{Exact preservation of native sequents}

\begin{theorem}[Witness-selected introduction]
\label{thm:det-witness-equivalence}
For every finite native context $\Gamma$ and object formula $A$,
\[
 \Gamma\vdash_{\mathcal C} A
 \quad\Longleftrightarrow\quad
 \Gamma\vdash_{\mathcal C_{\rm w}} A.
\]
Both translations are effective transformations of finite derivations.
The terms, formulas, contexts, and final sequent are unchanged.  In
particular the new rule has the original closure interpretation, not a
new progressive interpretation of atomic paired formulas.
\end{theorem}
\begin{proof}
Every application of DI$_{\rm w}$ is literally an old progress
application or an old case application.  The common rules are
unchanged.  Thus a new derivation is already an old derivation, proving
one inclusion without any translation of formulas.

For the reverse inclusion, induct on the old derivation, recursively
transforming its premise derivations first.  Common rules are repeated
unchanged.  Four introduction situations exhaust the remaining cases.

\paragraph{1. Old case, with compound matrix or nonliteral-zero witness.}
If $B$ is compound, or if $t\not\equiv_{\rm syn}0$, then
$\delta_{B,t}(z)=0$.  The transformed premise is exactly a premise of
DI$_{\rm w}$, so repeat the introduction.

\paragraph{2. Old case, with literal equation matrix and literal-zero
witness.}
Write the transformed premise as
\[
 \pi:\qquad\Gamma,E(0,0)\vdash E(0,Sz).
\]
Choose its eigenparameter $z$ fresh for all retained data.  Under the
new context $\Gamma,E(C,0)$, transport the hypothesis $E(C,0)$ along
the closed theorem $C=0$, using the fixed motive $E(w,0)$.  This gives
$E(0,0)$.  Lemma~\ref{lem:det-witness-cut} inserts this derivation
at the uses of the assumption $E(0,0)$ in $\pi$, giving
\[
                 \Gamma,E(C,0)\vdash E(0,Sz).
\]
Transport the resulting conclusion along $0=C$, now with motive
$E(w,Sz)$.  The conclusion is
\[
                 \Gamma,E(C,0)\vdash E(C,Sz).
\]
The term $C$ is closed and syntactically nonzero.  Consequently
DI$_{\rm w}$ with witness $C$ selects the case premise just obtained
and concludes the original formula $D_{x,y}E(x,y)$.  The conclusion's
syntax does not record which witness was used.

This argument changes the retained ordinary assumption by proved
transport and cut.  It does not apply a numerical substitution to
$\pi$, so it cannot accidentally change the selected mode of an
introduction inside that derivation.

\paragraph{3. Old progress, with literal-zero witness.}
Its matrix is a literal equation and its witness is the literal zero
term.  Therefore DI$_{\rm w}$ selects progress and repeats the old
inference using the recursively transformed premise.

\paragraph{4. Old progress, with a nonliteral-zero witness.}
Write its transformed premise as
\[
 \pi:\qquad\Gamma,E(t,z)\vdash E(t,Sz),
 \qquad t\not\equiv_{\rm syn}0.
\]
Choose a new variable $u$ absent from $\Gamma,t$ and the entire matrix.
Consider the equation template $\widetilde E(u,y):=E(t,y)$, whose
$u$ coordinate is unused.  DI$_{\rm w}$ with witness $0$ selects
progress and uses $\pi$ directly to give
\[
                    \Gamma\vdash D_{u,y}E(t,y).
\]
Paired elimination of this formula introduces the local closure
$Q_yE(t,y)$.  Its fresh existential parameter is absent from this
closure, because $u$ was unused.

Under $\Gamma,Q_yE(t,y),E(t,0)$, hypothesis closure and QE give
$E(t,Sz')$, for a newly chosen parameter $z'$ fresh for all retained
context and passive data.  Since $t$ is syntactically nonzero,
DI$_{\rm w}$ with witness $t$ selects cases and proves
\[
          \Gamma,Q_yE(t,y)\vdash D_{x,y}E(x,y).
\]
The paired elimination now discharges its local closure and gives
exactly the original conclusion.  The original premise proof $\pi$
has been inserted once.  Neither its witness nor its numerical
parameters have been substituted.

All fresh parameters used in these four constructions can be chosen
outside the finitely many variables of the transformed premises.
Renaming preserves the deterministic guard by
Lemma~\ref{lem:det-witness-cut}.  This completes the induction on the
old derivation and proves the second inclusion.
\end{proof}

\paragraph{Substitution and arithmetic repairs are consequences, not
premises.}
Numerical substitution can change a witness from a variable into the
literal constant $0$, and can also create a forbidden literal-zero
initializer in a restricted computation.  It is therefore not valid to
claim that arbitrary numerical substitution preserves every primitive
new rule instance.  Once Theorem~\ref{thm:det-witness-equivalence} is
proved, the inherited admissible substitution transformation may be
performed in the old calculus and its result translated back.  Likewise,
any separately established repair of the restricted computation axioms
transfers through the identity on native sequents.  The proof above used
neither assertion circularly.

\paragraph{Consequences and limits.}
The theorem transfers the inherited whole-domain bi-interpretation
with full first-order intuitionistic PRA, and its equiconsistency, to
$\mathcal C_{\rm w}$: compose the inherited proof translations with
the two identity-sequent transformations just proved.  All formula
round trips and numerical interpretation composites remain those of
the old calculus.  The new contribution is a restriction of primitive
proof choices, rather than a different numerical algebra or a new
meaning of intuitionistic logic.  No absolute minimum, lower-bound
optimality, or universal bound on the shortest certificates follows.
A concrete checker-size claim additionally requires the complete
implementation, its raw-input soundness argument, and measured source.

\paragraph{Attribution and the scope of the novelty search.}
The unused-witness progress construction is already present in the
inherited Choice-DI proofs, and the closed term $C$ and its direct zero
certificate are already present in the certified-zero draft.  Existing
progressive-binder proposals also use deterministic introduction, but
change the atomic closure interpretation and the corresponding
elimination rules.  The construction here combines the first two
inherited mechanisms to select the old case/progress rules by witness
syntax while preserving every old native sequent.  Searches of the
seventeen current extracted archives and their included text/code
references did not locate this exact restriction and simulation.
That bounded search is not a claim of literature-wide priority.

\paragraph{A general form of the normalization.}
The argument does not depend on the numerical value zero having a special
induction-theoretic role in the witness coordinate.  If $p$ is any fixed
closed term and there is a syntactically different closed term $q$ with
proofs of $p=q$ and $q=p$ in the common fragment, progress may instead be
selected exactly at witness $p$.  Repair a forbidden case witness $p$
using $q$, and introduce every auxiliary unused-witness progress formula
with witness $p$.  The rest of the simulation is identical.  Zero is
chosen here because its syntax is the cheapest recognizable witness and
its nonliteral representative $C$ has a direct computation certificate.
The induction coordinate itself still begins at zero; only the independent
existential witness coordinate is being normalized.


\section{The complete checker and its exact trust boundary}
\label{sec:executable}

This is the complete 555-byte source. The explanatory expansion in
\file{code/readable_checker.py} is independently executable. Neither
program imports a producer, a theorem database, a numerical evaluator,
or the other checker.

\begin{lstlisting}[language=Python]
import json
k = 0
i = lambda a, t=0j, d=0: [i(b, t, d + 5 - len(a)) for b in a] if a * 0 == [] else a if k <= a < d else d % 2 and i(t, d=3)

def c(a, g):
    global k
    if i(a) in g:
        return
    m, s, u, t = q()
    if m:
        c(m, g)
        a = m in ([s, u], not s and [t] * 3) and i(t, u) == a and i(t, s) or a
    h = m * 3 or (u and [[s and [s], u, t], i(t, [s, u, t], 3) if s else u]) or (a := i(*a, t, 2))
    k -= 1
    c(i(a, [k]), g + [i(h, (m or t == len(a) % 2 == 0) and k)])
q = json.load(open(0)).pop
c(q(), [])
print('ACCEPT')
\end{lstlisting}

\subsection{Certificate grammar and sorted interpretation}

Input is a JSON array. Its last element is the requested closed formula;
earlier four-field packets $[m,s,u,t]$ are consumed backwards. An unused
prefix is inert. Success means exit zero and exactly \texttt{ACCEPT}
followed by a newline. An exception rejects with a nonzero exit status.
The program does not use Python \texttt{assert}, so optimized Python
does not remove a logical check.

The arity representation is inherited~\cite{arity,certified}:
\begin{center}\small
\begin{tabular}{cll}
\toprule
Array arity & Numerical interpretation & Formula interpretation\\
\midrule
0 & Zero & Truth $0=0$\\
1 & Successor, with two dummy slots & Paired binder $D$\\
2 & Zero; annotations ignored & Equality\\
3 & Scalar iterator, one state slot & Closure $Q$ of its last field\\
At least 4 & Zero; children ignored & Truth; children ignored\\
\bottomrule
\end{tabular}
\end{center}
Canonical equations are $[s,t]$, paired formulas are $[B]$, and closures
are $[B,B,B]$. Numerical zero is scalar $0$, successor is unary, and an
iterator is ternary. Numerical roots have depth three; formula roots have
depth zero. Each traversal adds $5-\text{arity}$ to the depth. At a
numerical root, scalar positions $0,1,2$ are dummy zeros; successive
protected slots use pairs $3,4$, then $5,6$, and so on. Negative allocated
labels denote rigid parameters. In a successor its child has two dummy
zero slots; an iterator's counter and seed have one dummy slot, and the
update has its actual state slot. In $D$, the two slots are $y,x$; in
$Q$, the one slot is $y$.

The opener retains a scalar in $k\le a<d$. An escaping numerical
coordinate is filled by an image checked at numerical root depth three;
escaping formula padding erases to truth. The default image $0j$ is
intentionally not ordered against a real numerical coordinate. If it is
actually needed as a missing numerical image, the comparison fails,
just as it fails for the old \texttt{None}. Unused images need not be
evaluated. The list predicate \texttt{a * 0 == []} recognizes lists among
JSON values on which the expression succeeds; unsupported types reject.
The explicit-list variant is provided for a clearer spelling.

At ignored high arities the decoder is constant. The appropriate invariant
is semantic opening at meaningful sorted positions, conditional on a
successful traversal. A global raw-tree idempotence assertion is not
assumed. The inherited opening argument and its complete source are
identified in the provenance record; the new checker changes none of
its depth arithmetic or numerical substitution.

\subsection{The three operations and the shared continuation}

For a truthy major $m$, the checker proves $m$ in the retained context.
If $m=[s,u]$ and the target opening of $t$ is the current goal, it
reduces that goal by equality transport. If $s$ is falsy and
$m=[t,t,t]$, it performs closure specialization from its zero instance.
A failed endpoint match keeps the original goal. A unary paired major
exports its tripled matrix as a closure; other nonempty arrays export
constant truth. The recursive major proof is always required.

For a falsy major and a truthy initializer $u$, the packet generates
the equation
\[
 H=\begin{cases}
 I(0,u;v.t)=u,&s\text{ falsy},\\
 I(Ss,u;v.t)=t[I(s,u;v.t)/v],&s\text{ truthy}.
 \end{cases}
\]
Every selected generated equation is a nonempty two-element array.
Canonical forbidden zero seeds or zero predecessors are restored by
the certified-zero computation lemma. No separate arithmetic opcode
or semantic test that an initializer is positive is required.

For a falsy major and falsy initializer, the packet introduces a paired
goal, with $t$ as its witness. For canonical nonempty goals, unary
unpacking checks the paired shape. Once the witness is opened, progress is selected exactly by
\texttt{t == len(a) \% 2 == 0}. On canonical data this says that the
witness is literal zero and the opened matrix is an equation. Compound
matrices use their base instance. Larger harmless truth aliases retain
the inherited constant interpretation.

The continuation allocates a fresh negative name, opens the successor
goal, appends exactly one local hypothesis, and recurses with a new
context list. The context is not shared mutably with a major subproof.
Already checked goals and contexts do not mention names admitted only
after this allocation. JSON fractional numbers in a newly admitted
interval must be treated as distinct fresh names, not identified with
its integer endpoint; only finitely many occur in a finite trace.

\subsection{Why Boolean dispatch is legitimate}

The old assignment was a conditional on $m$, followed by a conditional
on $u$. In a successful new execution, a truthy permitted $m$ has a
truthy triple $m*3$. If no major is selected, a selected computation
has a nonempty equation array. Therefore the expression
\[
  m*3\ \texttt{or}\ (u\ \texttt{and}\ H)\ \texttt{or}\ \text{introduction}
\]
selects precisely the intended branch. Some old falsy marker aliases,
including null, fail at multiplication. Canonical arithmetic null
markers are replaced by $[]$. Both old markers are falsy and compare
unequal to an integer parity, so this replacement preserves the old
Choice-DI program's continuation. This is a restriction and normalization
of certificate data, not an additional logical inference.

\subsection{An exact simulation of every new accepting trace}

For a false-major packet executed by the witness checker, change its
marker to $0$ if its raw field $t$ equals $0$, and to $[]$ otherwise.
Leave every true major unchanged. Apply this fixed transformation to
the complete finite input stack. For a false major, the old guard
\[
             \texttt{len(a) \% 2 == m}
\]
then has exactly the new guard's value
\[
             \texttt{t == len(a) \% 2 == 0}.
\]
Indeed the first transformed marker compares equal only at parity zero,
and the second never compares equal to integer parity. This remains
true for arithmetic packets: their update occupies $t$, and the same
opening image is obtained. Both major branches, every recursive call,
and every new local assumption otherwise coincide. Replacing the unused
invalid-image default $0j$ by \texttt{None} cannot change a successful
trace, because either default raises if used as a numerical image.
Thus a new finite acceptance simulates an old certified-zero acceptance.

This argument includes noncanonical successful inputs. It avoids the
false assertion that every generated $H(b)$ is a uniformly true arithmetic
equation. In general only $H(\vec0)$ is known. Freshness permits substitution
of zero for all newly introduced names in a continuation proving a
conclusion independent of them, followed by cut against the known zero
instance. This substitution and cut take place in the simulated old
calculus. This is the inherited supply argument, not a new reflection rule.

For completeness, first use the theorem in
Section~\ref{sec:deterministic-witness} to normalize a finite named
derivation. Repair excluded computation leaves by the certified zero.
Serialize canonical terms and formulas with fresh names and use $[]$
for falsy packet markers. In a canonical computation all meaningful
parameters are already allocated, so opening at zero or at the new name
does not alter the supplied equation. The canonical introductions now
have exactly the mode selected by their witness. This constructs an
accepting finite certificate, assuming sufficient computation resources.

\section{Why this is full intuitionistic PRA}
\label{sec:bridge-summary}

The ordinary realization has actual equality, genuine falsity, the whole
numerical domain, the iterator equations, contextual literal-equation
induction, numerical coverage, and separation $(0=S0)\to\bot$.
The native formula interpretation is
\[
 (s=t)^\circ=(s=t),\quad
 (D_{x,y}B)^\circ=\exists x\bigl(B^\circ(x,0)\to\forall yB^\circ(x,y)\bigr),
 \quad (Q_yB)^\circ=B^\circ(0)\to\forall yB^\circ(y).
\]
The deterministic witness restriction does not change these meanings.
Its introductions are old case/progress instances, and its converse
simulation preserves the same end formulas and retained hypotheses.

For orientation, the inherited scalar selector is
$C(n,a,b)=I(n,a;v.b)$. Its zero and successor computations make it a
formula gate and a selector between arbitrary formulas after padding
their binder depths. The paired binder then defines implication,
conjunction, disjunction, existential quantification, and universal
quantification; $0=1$ represents falsity. Those definitions derive the
ordinary intuitionistic rules. Their explicit formulas, all introduction
and elimination arguments, and both logical composites are reproduced
in Appendix~\ref{sec:choice-introduction}. In particular the native
paired formula is interderivable with
\[
 \exists_*x\bigl(B(x,0)\Rightarrow_*\forall_*yB(x,y)\bigr).
\]
This round trip is essential; a progressive reinterpretation of atomic
$D$ would not suffice.

The numerical bridge follows the proved dependency order:
\begin{enumerate}
\item Equation induction and cases construct a scalar predecessor
      $P$ and prove $P0=0$, $P(Sn)=n$ internally.
\item These yield subtraction and a numerical test reflecting actual
      equality on the whole domain, rather than merely on standard values.
\item A characteristic-term construction transfers induction from
      literal equations to all fixed quantifier-free formulas.
\item Pairing and projections, established using that induction,
      compile each fixed general primitive recursion into scalar iteration.
\end{enumerate}
The full inherited proofs are reproduced in Appendices
\ref{sec:hybrid-bootstrap}, \ref{sec:eic-completion}, and
\ref{sec:iterator-coding}; they are not counted as imported checker code.

Conversely, iPRA defines every fixed scalar iterator by primitive recursion
that ignores its counter argument in the update. For each finite description,
the double translation satisfies the original defining equations.
Equation induction proves uniqueness of its values, at arbitrary internal
counters and with passive parameters fixed. Thus both numerical composites
are internally isomorphic to the original domain by $x=y$.
There is no restriction to a cut, retract, quotient, or standard numerals.
Both interpretations preserve falsity, using the explicit ordinary
separation sentence. Effective proof translations give equiconsistency;
no self-consistency theorem is claimed. Appendix~\ref{sec:ordinary-bridge}
spells out these composites.

\section{Reproducibility, evidence, and remaining scope}
\label{sec:evidence}

The new named proof normalizer is separate from the trusted checker.
It records the two actual mode repairs, preserves the final formula,
and then uses the unchanged attributed arity emitter. Its binder
convention uses globally distinct names; mathematical alpha-renaming
provides that convention for an arbitrary finite derivation.

The focused witness suite contains 18 positive and five negative cases.
It tests the three compact exact-list implementations and the readable
expansion, in ordinary and optimized Python: 184 complete CLI runs.
Two further mutation controls remove the literal-equation restriction;
the mutated program then accepts a forged compound-progress certificate,
whereas the delivered programs reject it. A zero-case test uses its base
assumption essentially, and a nonzero-progress test uses its induction
hypothesis with the witness occurring in the matrix. Their untranslated
certificates reject and the normalized certificates accept. These examples
exercise the two proof transformations, rather than only vacuous witnesses.

The unchanged-choice checker replays the 36 supplied certified-zero
cases after its marker translation, in both Python modes: 72 executions.
The independent initial baseline audit also recorded 216 executions of
the three upstream programs. Appendix~\ref{zcc:section} has six separate
symbolic congruence certificates and three side-condition rejection controls.
These suites overlap in mathematical content and are not counts of
distinct theorems. Finite checks support the implementation; the universal
soundness and equivalence conclusions rest on the written arguments.

From a fresh extraction, run:
\begin{lstlisting}[language=bash]
python3 reproduce.py
python3 code/checker.py < examples/zero_case_uses_base.json
python3 code/checker.py < examples/nonzero_progress_uses_IH.json
python3 code/checker.py < examples/closed_false.json
python3 build.py
\end{lstlisting}
The first two individual proof checks accept; the last rejects.
Reproduction verifies package hashes before rerunning tests. Rebuilding
or intentionally editing the source can change those hashes, particularly
PDF metadata, so use a working copy if the original manifest is to remain
an identity record. The build produces a standalone LaTeX source and
compiles without shell escape. The only build dependency outside Python's
standard library is a suitable TeX installation.

The report and code establish a small positive result under a fixed,
reproducible comparison. They do not establish the significant breakthrough
requested. Removing an independent certificate choice can simplify a rule
inventory while increasing another implementation cost: here its guard
adds five bytes. The smaller overall programs should therefore not be
advertised as evidence of a radically simpler numerical foundation.

\clearpage
\appendix
\section*{Attributed complete bridge proofs}
\addcontentsline{toc}{section}{Attributed complete bridge proofs}
The following five sections reproduce the corresponding source fragments
from the supplied certified-zero package, byte-for-byte. They document the
inherited theory to which the new identity-sequent theorem applies. Their
historical constructions are not new contributions of this report.
% New mathematical result: complementary introduction choices.
% Requires amsmath/amssymb/amsthm and theorem, lemma, proposition environments.
% The inherited arithmetic-to-iPRA bridge must be supplied separately.

\section{Complementary introductions with one local assumption}
\label{sec:choice-introduction}

The old paired introduction retains both the zero instance and an
induction instance when its matrix is an equation. Both together are
unnecessary. A certificate may choose either one, provided the induction
choice is restricted to a literal equation. This changes neither the
formula language nor the meaning of a formula. The omitted combination
is recovered by a short nested introduction--elimination construction.

\subsection{Complete native syntax and rule inventory}

The scalar numerical language is
\[
 t::=x\mid0\mid St\mid I(n,a;u.h).
\]
Only $h$ binds the state variable $u$; $n$ and $a$ do not.
The update may contain fixed numerical parameters. The computation
axioms in every retained context are
\[
 I(0,a;u.h)=a,\qquad
 I(Sn,a;u.h)=h[I(n,a;u.h)/u].
 \tag{I0, IS}\label{eq:choice-computations}
\]
The argument also applies unchanged to the full fixed-description PR
language with its projection, composition, and recursion equations.
No numerical variable ranges over function descriptions, and no uniform
evaluator is postulated.

Native formulas and reusable context clauses are
\[
 A::=(s=t)\mid\mathsf D_{x,y}A,\qquad
 \langle A;y.B(y)\rangle,\quad y\notin\operatorname{FV}(A).
\]
Bound variables are renamed before capture-free substitution.
The two variables $x,y$ in a paired formula are distinct.
An ordinary assumption $F$ is the constant clause
$\langle H;v.F\rangle$, where $v$ is absent from $F$ and $H$ is the
fixed computation theorem $I(0,0;u.0)=0$. Thus ordinary hypothesis
use is a clause Use preceded by this computation proof.

There are computation axioms and exactly the following inference
patterns:
\[
\frac{\Gamma\vdash A}{\Gamma\vdash B(t)}
\quad\mathsf{Use}
\quad(\langle A;y.B(y)\rangle\in\Gamma),
\]
\[
\frac{\Gamma\vdash\mathsf D_{x,y}B(x,y)\qquad
      \Gamma,\langle B(a,0);y.B(a,y)\rangle\vdash F}
     {\Gamma\vdash F}
\quad\mathsf{DE},
\qquad
\frac{\Gamma\vdash s=t\qquad\Gamma\vdash F(s)}
     {\Gamma\vdash F(t)}
\quad\mathsf{Tr},
\]
\[
\frac{\Gamma,B(t,\eta_\varepsilon(z))\vdash B(t,Sz)}
     {\Gamma\vdash\mathsf D_{x,y}B(x,y)}
\quad\mathsf{DI}_{\varepsilon},
\qquad
\eta_0(z)=0,\quad\eta_1(z)=z.
\tag{Choice DI}\label{eq:choice-di}
\]
Here $\varepsilon$ is a certificate choice in $\{0,1\}$, and
$\varepsilon=1$ is permitted only when $B$ has a literal equation
root. An equation beneath a binder is not a literal equation.
Call $\varepsilon=0$ the \emph{case choice} and $\varepsilon=1$
the \emph{progress choice}. Both retain exactly one ordinary
formula assumption. Compound matrices have only the case choice.

In DI the eigenparameter $z$ is fresh for $\Gamma$, the witness
$t$, and every fixed parameter of $B$. In DE the eigenparameter
$a$ is fresh for $\Gamma$, the conclusion $F$, and every fixed
parameter of $B$. In transport, $F(v)$ is any fixed native
formula template. Displayed local assumptions and clauses are
discharged. There are no primitive reflexivity, coverage,
separation, ordinary logical, or additional induction rules.

We call this calculus $\mathcal C$. The original comparison
calculus $\mathcal N$ has identical syntax, computation,
Use, DE, and transport, but its introduction is
\[
\frac{\Gamma,B(t,0),[B(t,z)]_{\mathrm{at}}\vdash B(t,Sz)}
     {\Gamma\vdash\mathsf D_{x,y}B(x,y)}.
\tag{Old DI}\label{eq:choice-old-di}
\]
The bracket supplies $B(t,z)$ exactly at a literal equation.

\subsection{Exact equivalence and the unused-witness construction}

\begin{lemma}[Structural admissibility]
\label{lem:choice-structural}
Weakening, exchange, contraction, renaming, ordinary-assumption cut,
and capture-free numerical substitution are admissible in $\mathcal C$.
\end{lemma}
\begin{proof}
Induct on derivations. Rename local eigenparameters before inserting
proofs or terms. Numerical substitution preserves a formula's root
constructor, so an allowed progress choice remains allowed.
The selected choice does not change under substitution or cut.
For ordinary-assumption cut, replace each use of its constant clause
by the supplied proof. Computation schemas are closed under
capture-free numerical substitution. These facts verify every case.
\end{proof}

\begin{theorem}[Identity-language equivalence]
\label{thm:choice-equivalence}
For every retained clause context $\Gamma$ and native formula $F$,
\[
 \Gamma\vdash_{\mathcal C}F
 \quad\Longleftrightarrow\quad
 \Gamma\vdash_{\mathcal N}F.
\]
The transformations are effective on finite derivations and leave
all numerical and formula syntax unchanged.
\end{theorem}
\begin{proof}
Each new introduction is an instance of the old one after weakening
by the omitted local assumption. For the case choice at an equation,
add the induction instance. For the progress choice, add the zero
instance. Compound case introduction is already the old rule.
Thus every new derivation translates to an old one.

For the reverse direction, transform an old derivation recursively.
Only its atomic DI case needs a construction. Suppose its matrix is
the equation $E(x,y)$ and its transformed premise is
\[
 \pi:\quad\Gamma,E(t,0),E(t,z)\vdash E(t,Sz).
\]
The outer new inference will be case DI with witness $t$. It is
therefore enough to prove $E(t,Sz)$ under
$\Gamma,E(t,0)$.

Rename the eigenvariable in $\pi$ to a fresh $v$. Choose a new
variable $u$ absent from $E(t,y)$ and all retained data. Under
$\Gamma,E(t,0)$, progress DI with witness zero introduces
\[
             \mathsf D_{u,y}E(t,y)
\]
from exactly the renamed premise $\pi[v/z]$.
The bound coordinate $u$ is deliberately unused.

Apply DE to this paired formula. Because $u$ is unused, the opened
clause is exactly
\[
               \langle E(t,0);y.E(t,y)\rangle;
\]
it contains no new existential witness. The retained ordinary
assumption $E(t,0)$ supplies its antecedent. Use at $Sz$ gives
$E(t,Sz)$, and DE closes. Finally the outer case DI introduces
the original $\mathsf D_{x,y}E(x,y)$ with witness $t$.
The original premise proof $\pi$ has been inserted once.

An old compound introduction already is the new case choice.
All other rule cases are unchanged. This completes the induction
on the original derivation.
\end{proof}

The unused inner coordinate is essential. Introducing
$\mathsf D_{x,y}E(x,y)$ internally and then eliminating it would
give an arbitrary fresh witness, not necessarily the original
$t$. The construction instead keeps $t$ as a fixed parameter
of $E(t,y)$ and binds only the unused coordinate $u$.

\begin{proposition}[Linear bound on inference nodes]
\label{prop:choice-node-bound}
Count each computation axiom, Use, DI, DE, and transport as one
inference node, expanding ordinary-assumption use as a computation
of $H$ followed by Use. If an old derivation has $N$ nodes and
$r$ atomic DI nodes, the displayed transformation has exactly
\[
                   N+5r\leq 6N
\]
nodes.
\end{proposition}
\begin{proof}
Each atomic old DI node is replaced by six nodes: outer case DI,
inner progress DI, DE, Use of the opened clause, Use of the
retained base assumption, and the computation of $H$ needed for
that last ordinary-assumption use. Its old premise subtree
occurs once. All other inference nodes are copied once.
\end{proof}

This counts inference-tree nodes. It is not a bound on serialized
JSON bytes, formula sizes, checker runtime, or the number of
consumed packet cells when hypothesis closure is implicit.
In the formula-context Q/D presentation, the corresponding
replacement has outer case DI, inner progress DI, DE, QE,
and two Hyp nodes, giving the same six-node count.

\subsection{Induction and numerical cases separate cleanly}

\begin{lemma}[Contextual equation induction]
\label{lem:choice-ei}
The progress choice, DE, and Use derive
\[
\frac{\Gamma\vdash E(0)\qquad
      \Gamma,E(z)\vdash E(Sz)}
     {\Gamma\vdash E(t)}
\tag{EI}\label{eq:choice-ei}
\]
for every equation $E$, with the usual freshness condition.
\end{lemma}
\begin{proof}
Choose an unused witness coordinate $u$. Progress DI introduces
$\mathsf D_{u,y}E(y)$ from the given step. DE opens the clause
$\langle E(0);y.E(y)\rangle$. Supply its antecedent using the given
base proof and use the clause at $t$. Close DE.
\end{proof}

\begin{lemma}[Cases, equality, and binder congruence]
\label{lem:choice-preinduction}
The case choice derives, for arbitrary native formulas,
\[
\frac{\Gamma\vdash F(0)\qquad\Gamma\vdash F(Sz)}
     {\Gamma\vdash F(t)}
\tag{NC}\label{eq:choice-nc}
\]
with the usual freshness condition. Equality reflexivity,
symmetry, transitivity, numerical congruence, and uniform
matrix interderivability under $\mathsf D$ are derivable
without using the progress choice.
\end{lemma}
\begin{proof}
For cases, use case DI with an unused witness coordinate and
matrix $F(y)$, ignoring its retained base in the successor
premise. DE opens $\langle F(0);y.F(y)\rangle$; the given zero
proof and Use at $t$ finish.

For reflexivity put $r=I(0,t;u.0)$. Use the computation equation
$r=t$ as both transport premises with motive $v=t$ to obtain
$t=t$. The usual explicit equality motives give symmetry,
transitivity, and numerical congruence.

For binder congruence, open the old paired formula at a fresh
witness $a$ and introduce the new one by case DI at $a$.
Convert its retained base to the old base, use the opened
old clause at the successor, and convert back. Ordinary cut
closes the construction, and the reverse conversion gives
the converse. No induction is used.
\end{proof}

The primitive progress choice is still restricted to equations.
Allowing it at arbitrary compound matrices would, by the
argument for \eqref{eq:choice-ei}, supply induction for
arbitrary formulas and change the intended theory.

\subsection{The entire ordinary logical interface}
\label{subsec:choice-logic}

Here is a predecessor-free reconstruction of the ordinary logical
operations. It makes explicit that the case choice supplies the
logical interface independently of progress induction.

Let $C(v,a,b):=I(v,a;u.b)$ with $u$ unused. Computation gives
$C(0,a,b)=a$ and $C(Sv,a,b)=b$. Define the gate on fixed finite
formula syntax by
\[
 G_v(s=t):=(C(v,0,s)=C(v,0,t)),\qquad
 G_v(\mathsf D_{x,y}B):=\mathsf D_{x,y}G_vB,
\]
renaming bound variables away from the control. Equality,
binder congruence, and external structural induction give
\[
 \vdash G_0A,\qquad G_{Sv}A\dashv\vdash A,\qquad
 A\dashv\vdash\mathsf D_{u,y}G_yA
 \quad(u,y\notin\operatorname{FV}(A)).
\tag{Gate and padding}\label{eq:choice-gate}
\]
For padding, case DI proves the forward direction. DE, the
zero-gate theorem, and Use at $1$ prove the reverse direction.

Define a selector between two arbitrary formulas by
\[
 M_v(s=t,u=w):=(C(v,s,u)=C(v,t,w))
\]
and, with aligned fresh paired variables, by
\[
\begin{aligned}
 M_v(\mathsf D_{x,y}A,\mathsf D_{x,y}B)
   &:=\mathsf D_{x,y}M_v(A,B),\\
 M_v(E,\mathsf D_{x,y}B)
   &:=\mathsf D_{x,y}M_v(G_yE,B),\\
 M_v(\mathsf D_{x,y}A,E)
   &:=\mathsf D_{x,y}M_v(A,G_yE),
\end{aligned}
\]
where $E$ is an equation. Every call decreases maximum input
binder depth. Padding and congruence prove
$M_0(A,B)\dashv\vdash A$ and
$M_{Sv}(A,B)\dashv\vdash B$.

Put $\delta(v)=C(v,1,0)$ and $p(z)=I(z,0;u.\delta(u))$.
Only the fixed computations $p(1)=1$ and $p(2)=0$ are needed.
The full logical macro inventory is
\[
\begin{aligned}
 A\Rightarrow_*B&:=\mathsf D_{u,z}M_z(A,B),\\
 A\wedge_*B&:=\mathsf D_{u,z}G_zM_{p(z)}(A,B),\\
 \exists_*x\,A(x)&:=\mathsf D_{x,y}G_yA(x),\\
 A\vee_*B&:=\exists_*x\,M_x(A,B),\\
 \forall_*x\,A(x)&:=
        A(0)\wedge_*\mathsf D_{u,z}G_zA(z),\\
 \bot_*&:=(0=1).
\end{aligned}
\tag{Logical macros}\label{eq:choice-macros}
\]
Newly bound variables are fresh; $u$ is unused.

Implication introduction obtains its assumption from the zero
selector and supplies the consequent at successors, then uses
case DI. Elimination opens its clause and uses it at $1$.
Existential introduction uses its witness and the successor
gate; elimination opens at a fresh witness, supplies the
zero-gate theorem, and uses the clause at $1$.

Disjunction introduction uses witnesses $0$ and $1$.
For elimination, discharge the two branch proofs as
$A\Rightarrow_*F$ and $B\Rightarrow_*F$. NC applied to
$M_z(A,B)\Rightarrow_*F$ gives this implication at every
control: the two endpoint cases follow from the branch
implications. Existential elimination provides $M_a(A,B)$
and the implication gives $F$.

Given $A$ and $B$, selector endpoints and NC give $M_v(A,B)$
at an arbitrary control. Substitute $v=p(Sz)$ and use the
successor gate for conjunction introduction. For projections,
open its zero-gated clause and specialize at $1$ and $2$;
the two fixed computations give $B$ and $A$, respectively.

These rules prove existential numerical coverage:
\[
     (x=0)\vee_*\exists_*v\,(x=Sv).
\tag{Coverage}\label{eq:choice-coverage}
\]
Use NC. Zero takes the left injection of reflexivity, and a
successor takes the right injection with the evident witness.
Coverage gives equality cases even when the retained context
contains $x$; its elimination adds $x=0$ or $x=Sv$ rather
than silently changing that retained context.

Universal introduction follows from a fresh-parameter proof
by substituting zero and successors and forming the two
conjuncts. For elimination, take its $A(0)$ conjunct and
open the paired conjunct. Its base is a zero-gate theorem.
Coverage of the target chooses the zero instance from
$A(0)$ or a successor instance from the opened clause
and gate removal.

Finally, under $0=1$, computation and transport give
$s=C(0,s,t)=C(1,s,t)=t$ for arbitrary numerical terms.
Every equation follows, and external induction on formula
syntax with case DI gives every native formula. This
supplies falsity elimination. No excluded middle for
quantified formulas has entered the construction.

\subsection{Whole-domain realizations and both formula composites}

For a standard first-order reading, let $T_{\rm eq}$ have
intuitionistic equality logic, the same numerical signature
and computation equations, numerical coverage, separation
$(0=1)\to\bot$, and contextual equation induction.
Its interpretation of the native syntax is
\[
 (s=t)^\circ=(s=t),\qquad
 (\mathsf D_{x,y}B)^\circ
   =\exists x\bigl(B^\circ(x,0)\to\forall y B^\circ(x,y)\bigr),
\]
\[
 \langle A;y.B\rangle^\circ
       =A^\circ\to\forall y B^\circ.
\]
The domain is the entire numerical domain and interpreted
equality is actual equality.

Case DI is sound in this realization by numerical coverage.
Progress DI is sound by contextual equation induction under
the conditional base. Use, DE, transport, and computation
have their displayed ordinary meanings. Conversely, the
native macros derive every ordinary logical rule,
\eqref{eq:choice-ei} derives equation induction, and
\eqref{eq:choice-coverage} derives coverage. Separation
translates to the identity implication
$(0=1)\Rightarrow_*(0=1)$.

For the essential paired formula composite,
\[
 \mathsf D_{x,y}B
 \dashv\vdash
 \exists_*x\bigl(B(x,0)\Rightarrow_*\forall_*y\,B(x,y)\bigr).
\tag{Paired composite}\label{eq:choice-paired-composite}
\]
Forwards, DE opens its clause at a fresh $a$; under $B(a,0)$
that clause supplies $B(a,y)$ at a fresh arbitrary $y$.
Introduce the universal, implication, and existential
macros and close DE. Backwards, eliminate the existential
at a fresh $a$. Introduce the paired formula with witness
$a$ by case DI: the retained base, supplied implication,
and universal elimination at $Sz$ give its successor premise.

The ordinary interpretations of the six logical macros
have their ordinary round trips. For conjunction, its true
base and the instances $1,2$ recover both operands;
the converse uses selector cases. The universal has all
successor instances in its paired conjunct and zero in
its first conjunct; coverage gives exactly the ordinary
universal. Implication, disjunction, and existential
quantification use their selector or gate endpoints.
Falsity uses the explicit ordinary separation bridge.
External induction on finite formulas, together with
\eqref{eq:choice-paired-composite}, gives both formula
composites. Clauses correspond in both directions to
$A\Rightarrow_*\forall_*yB$ by the derived logical rules.

The numerical composites are literal identities with
internally provable graph $x=y$. No standard-numeral
restriction, retract, quotient, or equality reflection
assumption is being inserted at this stage. The separately
supplied arithmetic bridge identifies $T_{\rm eq}$ with
full iPRA; the present change preserves that bridge
unchanged because Theorem~\ref{thm:choice-equivalence}
preserves every native sequent.

\subsection{What the theorem does and does not claim}

The new inference interface retains one ordinary assumption
at every paired introduction. It preserves the entire
consequence relation, not just equations or closed
theorems. Consequently the established full-iPRA
interpretations and equiconsistency transfer unchanged.

This result is different from retaining only progress
introduction and strengthening closure elimination.
Here the original elimination rule is retained, and
the case and progress introductions coexist. It is also
different from allowing arbitrary compound progress,
which would supply stronger induction.

A complete-checker comparison must still count the
certificate choice, its equation-root restriction,
substitution, context handling, computation, and the
public input/output interface. The mathematical theorem
alone is not a byte-count or minimality theorem.
The inference-node bound is distinct from source-size
measurements and from serialized certificate size.
The proof is a written universal proof; finite regression
tests do not constitute a Lean/Rocq formalization of it.

% Adapted with attribution from the pinned unparenthesized one-hole archive.
\section{Internal scalar predecessor bootstrap}
\label{sec:hybrid-bootstrap}

This scalar bootstrap is inherited from the supplied scalar development~\cite{erasure} and is reproduced to keep the scalar checker self-contained. It is not a new contribution of this paper. The proof uses a factor map fixing odd states and a saturated-iteration construction of predecessor.

The available arithmetic terms are $0$, $S$, and
\[
 I(n,a;u.h),\qquad
 I(0,a;u.h)=a,\qquad
 I(Sn,a;u.h)=h[I(n,a;u.h)/u].
\]
The update binds only $u$.  All other variables in it are fixed numerical
parameters.  We use equality reasoning, the already reconstructed
zero/successor cases, and induction on a single literal equation.
Every induction below keeps its other parameters fixed.  Predecessor,
successor injectivity, pairing, and induction on arbitrary
quantifier-free formulas are not assumptions of the argument.

\subsection{Elementary terms and two finite cycles}

Write $1=S0$ and use larger numerals as successor abbreviations.  Define
\[
 C(n,a,b)=I(n,a;u.b),\qquad
 Z(n)=C(n,1,0),\qquad
 a+b=I(b,a;u.Su).
\]
The update variable of $C$ is absent from its body.  Consequently
\[
 C(0,a,b)=a,\qquad C(Sn,a,b)=b.
\]
Numerical cases show that $Z(n)$ is $0$ or $1$.
The usual commutative-monoid equations for addition follow by literal
equation induction.  For example, $0+b=b$ and $Sa+b=S(a+b)$ are
proved by induction on $b$; these yield commutativity by a further
equation induction, and associativity follows by induction on the last
summand.  We use these equations to reorder finite sums.  The notation
$2a$ means $a+a$ and $3a$ means $a+a+a$; no multiplication operation
is needed for this bootstrap.

Put
\[
 N(n)=I(n,0;u.Z(u)),\qquad
 j(u)=Z(u)+2N(u),\qquad
 M(n)=I(n,0;u.j(u)).
\]
The range of $N$ is contained in $\{0,1\}$: its zero value is $0$,
and its successor value is a value of $Z$.  This uses numerical cases
on its counter, not induction on a disjunction.  Hence
\begin{equation}
 N(n)+Z(N(n))=1,\qquad
 Z(Z(N(n)))=N(n),\qquad N(n+2)=N(n).
 \label{eq:hybrid-parity}
\end{equation}
Here the last equation follows from two successor computations and
the preceding equation.  Literal-equation inductions give
\begin{equation}
 N(2n)=0,\qquad N(2n+1)=1,\qquad N(3n)=N(n).
 \label{eq:hybrid-parity-scaling}
\end{equation}
For the first equation the successor step adds two to the argument of
$N$; the second follows by one computation.  In the third, increasing
$n$ adds three to the left argument and therefore applies $Z$ once to
both sides of the induction equation.

Finite computations now give
\[
 j(0)=1,\qquad j(1)=2,\qquad j(2)=0.
\]
The proof below needs just these three values and
$M(Sn)=j(M(n))$; it does not assume a division algorithm or a theorem
characterizing $M$ as remainder modulo three.

\subsection{A ramp with a global scaling identity}

Define
\[
 G(n)=I(n,0;u.u+M(u)+1).
\]
The central invariant is
\begin{equation}
 M(G(n))=N(n).
 \label{eq:hybrid-ramp-invariant}
\end{equation}
Its zero case is immediate.  For its successor step, retain the
single induction equation and split the Boolean value $N(n)$.
If it is $0$, the next $G$ state is $G(n)+1$; its $M$ value is
$j(0)=1$, equal to $N(Sn)$.  If it is $1$, the next $G$ state
is $G(n)+2$; two computations of $M$ change its value from $1$
to $2$ and then to $0$, again equal to $N(Sn)$.  This proves
the invariant by literal-equation induction.

It follows that
\[
 G(Sn)=G(n)+N(n)+1,\qquad G(n+2)=G(n)+3.
\]
Equation induction therefore yields
\begin{equation}
 G(2n)=3n.
 \label{eq:hybrid-ramp-double}
\end{equation}
A separate literal-equation induction proves
\begin{equation}
 G(3n)=3G(n)+N(n).
 \label{eq:hybrid-ramp-triple}
\end{equation}
For completeness, its successor calculation is
\begin{align*}
 G(3Sn)
 &=G(3n)+N(3n)+4
   =3G(n)+2N(n)+4,\\
 3G(Sn)+N(Sn)
 &=3G(n)+3N(n)+3+Z(N(n)).
\end{align*}
The expressions agree in both Boolean cases for $N(n)$.
The base is $0=0$.  In particular,
\eqref{eq:hybrid-ramp-triple} was proved internally; it was not
inferred from an intended floor-function interpretation of $G$.

Set
\[
 H(u)=C(N(u),G(u),u).
\]
The crucial choice is to retain $u$ when its parity is $1$.
The preceding equations give
\begin{equation}
 H(2u)=3u,\qquad H(1)=1,\qquad H(3u)=3H(u).
 \label{eq:hybrid-factor-map}
\end{equation}
To check the last identity, use $N(3u)=N(u)$ and split the Boolean
$N(u)$.  At $0$, both sides are $3G(u)$ by
\eqref{eq:hybrid-ramp-triple}; at $1$, both are $3u$.
This argument does not require that an arbitrary number has already
been represented as $2v$ or $2v+1$.

\subsection{A comparator from scalar iteration}

For every fixed update $h$, equation induction on $n$, keeping $m,a$
and all parameters fixed, proves the semigroup identity
\[
 I(m+n,a;h)=I(n,I(m,a;h);h).
\]
The base is the zero computation and the step applies $h$ to the
induction equation. This identity will also be used in the later
pairing construction.


For a fixed update $h$, equation induction on $y$, with its initializer
$a$ and all parameters of $h$ fixed, proves the rotation identity
\begin{equation}
 I(Sy,a;h)=I(y,h(a);h).
 \label{eq:hybrid-rotation}
\end{equation}
Both sides have base $h(a)$, and their successor equations apply the
same $h$ to the previous equality.  Here and below $I(y,a;h)$ is
an abbreviation for one fixed finite scalar term, not quantification
over functions.

By \eqref{eq:hybrid-factor-map}, another equation induction gives
\begin{equation}
 I(y,3a;H)=3I(y,a;H).
 \label{eq:hybrid-iterator-scaling}
\end{equation}
Its base is $3a=3a$.  In the step the fixed-parameter induction
equation rewrites the new left side as
\[
 H(3I(y,a;H))=3H(I(y,a;H)).
\]
The initializer parameter stays $a$ throughout this induction.
Equation induction also gives $I(y,1;H)=1$, since $H(1)=1$.

Define
\[
 E(x)=I(x,1;u.2u),\qquad
 L(x,y)=N\bigl(I(y,E(x);H)\bigr).
\]
Then the following entire recursion specification is provable:
\begin{equation}
 L(0,y)=1,\qquad L(Sx,0)=0,\qquad L(Sx,Sy)=L(x,y).
 \label{eq:hybrid-comparator}
\end{equation}
The first equation uses $E(0)=1$ and the fixed-point identity for
$H$.  The second uses $E(Sx)=2E(x)$ and $N(2a)=0$.
For the third,
\begin{align*}
 L(Sx,Sy)
 &=N\bigl(I(Sy,2E(x);H)\bigr)\\
 &=N\bigl(I(y,H(2E(x));H)\bigr)\\
 &=N\bigl(I(y,3E(x);H)\bigr)\\
 &=N\bigl(3I(y,E(x);H)\bigr)\\
 &=N\bigl(I(y,E(x);H)\bigr)=L(x,y).
\end{align*}
The successive steps use rotation, $H(2a)=3a$, iterator scaling,
and $N(3a)=N(a)$.  Thus $L$ has the comparison equations at all
arguments internally.  Its standard-model meaning as the indicator
of $x\le y$ is motivation, not a premise of the following argument.

\subsection{Saturated iteration gives predecessor}

We now need only the three equations
\eqref{eq:hybrid-comparator}.  Define
\[
 B_m(u)=C(L(Su,m),u,Su),\qquad
 W(t,m)=I(t,0;u.B_m(u)).
\]
The subscript in $B_m$ records a fixed parameter; $B_m$ is an open
term abbreviation, not a function variable.

Numerical cases on $e$ prove the selector equation
\[
 C(e,Sa,Sb)=S C(e,a,b).
\]
Equivalently one may use equation induction on $e$, whose induction
hypothesis is unused.  Therefore
\begin{align}
 B_{Sm}(Su)
 &=C(L(SSu,Sm),Su,SSu)\notag\\
 &=C(L(Su,m),Su,SSu)=S B_m(u),
 \label{eq:hybrid-conjugacy}\\
 B_{Sm}(0)&=C(L(S0,Sm),0,S0)=C(L(0,m),0,S0)=S0.
 \label{eq:hybrid-conjugacy-base}
\end{align}

Literal-equation induction on $t$, with $m$ fixed, proves
\begin{equation}
 W(St,Sm)=S W(t,m).
 \label{eq:hybrid-saturated-shift}
\end{equation}
The base is \eqref{eq:hybrid-conjugacy-base}.  The successor step is
\begin{align*}
 W(SSt,Sm)
 &=B_{Sm}(W(St,Sm))\\
 &=B_{Sm}(S W(t,m))\\
 &=S B_m(W(t,m))=S W(St,m).
\end{align*}
Only the second line uses the current induction hypothesis, at the
same fixed $m$.

Two subsequent equation inductions give
\begin{equation}
 L(Sn,n)=0,\qquad W(n,n)=n.
 \label{eq:hybrid-diagonals}
\end{equation}
The first uses the comparator recurrence.  The second uses
\eqref{eq:hybrid-saturated-shift} in its successor step.  The
substitution $t=m=n$ is made into that already proved universally
parameterized equation; no parameter of an active induction hypothesis
is changed.  Consequently
\begin{equation}
 W(Sn,n)=B_n(n)=C(L(Sn,n),n,Sn)=n.
 \label{eq:hybrid-extra-step}
\end{equation}

Finally define the finite scalar term
\begin{equation}
 P(m)=I\bigl(m,0;u.C(L(SSu,m),u,Su)\bigr).
 \label{eq:hybrid-predecessor-term}
\end{equation}
Its zero equation is immediate.  To justify its successor equation
without assuming a separate extensionality rule for iterator bodies,
put
\[
 V(t,m)=I\bigl(t,0;u.C(L(SSu,Sm),u,Su)\bigr).
\]
The comparator recurrence identifies this update pointwise with
$B_m(u)$.  Equation induction on $t$, with $m$ fixed, gives
$V(t,m)=W(t,m)$: their base values are $0$, and the successor step
uses the pointwise update equation and then the induction equation.
Now \eqref{eq:hybrid-extra-step} yields
\[
 \boxed{\qquad P(0)=0,\qquad P(Sm)=V(Sm,m)=W(Sm,m)=m.\qquad}
\]

\subsection{Dependency boundary}

All terms in this section are built from $0,S,I$ by finite substitution.
All inductions have a single equation as their invariant.  Numerical
cases and equality transport are the already reconstructed logical
resources.  Successor injectivity may now be obtained by applying $P$
to an assumed equation $Sx=Sy$.  Equality tests, reflection of
quantifier-free formulas, pairing, and the recovery of general
primitive recursion can then be developed in the established order.
They were not used in order to manufacture predecessor.

This argument is a written internal derivation in the stated theory.
The finite implementation checks do not constitute a proof-assistant
formalization of this entire metatheorem.

% Standalone section.  Requires only amsmath and amssymb in the parent document.
% No theorem environments, macros, or bibliography keys are assumed.
% All arithmetic claims in this section are proved here.

\section{From literal-equation induction to full intuitionistic PRA}
\label{sec:eic-completion}\label{app:arithmetic}

This section supplies an unconditional completion theorem. It is adapted
with attribution from the current forward-left draft~\cite{forward}. In the
scalar application used here, every function occurring in the proof is
already a scalar term: predecessor is supplied by
Appendix~\ref{sec:hybrid-bootstrap}, the selector is
$C(n,a,b)=I(n,a;u.b)$, addition is $x+y=I(y,x;u.Su)$, and
truncated subtraction is $d(x,y)=I(y,x;u.Pu)$. All other functions below
are finite compositions of these. Thus this proof does not presuppose
that arbitrary primitive recursion has already been compiled. The QFI
conclusion first applies to the scalar language; the general PR compiler
comes afterward.  Literal-equation
induction, together with the primitive-recursive defining equations,
predecessor coverage, and separation of $0$ and $1$, proves induction for every
quantifier-free formula.  In particular, the proof does not assume a faithful
numerical equality test, stability or decidability of equality, successor
uniqueness, or any commutation law for an iterator.  We construct the equality
test and prove its reflection property for arbitrary elements of the
numerical domain.

\subsection{The exact hypotheses}

Work in ordinary intuitionistic first-order logic with equality and genuine
falsity $\bot$.  The language has one numerical sort and a symbol for each
fixed, finite primitive-recursive description.  Distinct descriptions are
not identified merely because they compute the same standard function.
Let $\mathsf{PRE}$ contain the defining equations for zero, successor,
projections, composition, and primitive recursion.  In particular, choose
descriptions $P$ and $C$ with equations
\begin{align}
 P(0)&=0, & P(Sx)&=x, \label{eq:eic-predecessor}\\
 C(0,a,b)&=a, & C(Sx,a,b)&=b. \label{eq:eic-selector}
\end{align}
Write $1=S0$.  In addition to $\mathsf{PRE}$, assume
\begin{align}
 \mathsf{Cov}:&\qquad
       \forall x\,\bigl(x=0\lor x=S(Px)\bigr),
       \label{eq:eic-coverage}\\
 \mathsf{Sep}:&\qquad 0=1\longrightarrow\bot,
       \label{eq:eic-separation}
\end{align}
and induction for each \emph{literal equation} $E(n,\vec p)$:
\begin{equation}
 \forall\vec p\,
 \left[
   \left(E(0,\vec p)\land
       \forall n\bigl(E(n,\vec p)\longrightarrow E(Sn,\vec p)\bigr)
   \right)
   \longrightarrow\forall n\,E(n,\vec p)
 \right].
 \label{eq:eic-induction}
\end{equation}
Here the schematic expression $E$ is a single equality of numerical
terms; the displayed universal quantifiers range only over numbers.
Every description and every instance of the schema is finite external
syntax.

The axiom formulation \eqref{eq:eic-induction} and the usual contextual
induction rule are interchangeable.  For example, if
\[
 \Gamma\vdash E(0,\vec p),
 \qquad
 \Gamma,E(z,\vec p)\vdash E(Sz,\vec p),
\]
where $z$ is absent from $\Gamma$ and the passive parameters, discharge the
second premise's equation and universally introduce $z$.  The axiom then
gives $\Gamma\vdash\forall n E(n,\vec p)$, which can be instantiated at any
target term.  Conversely, assume the antecedent of the axiom and apply the
contextual rule with fresh variables.  This uses the ordinary intuitionistic
logical rules and permits arbitrary, including quantified, formulas in
$\Gamma$.

\subsection{Constructor separation and two elementary tests}

Define the sign description and the countdown description by
\begin{equation}
 \operatorname{sg}(u)=C(u,0,1),
 \qquad
 K(u)=C(u,1,Pu).
 \label{eq:eic-sign-countdown}
\end{equation}
Their defining equations imply
\begin{align}
 \operatorname{sg}(0)&=0,
 &\operatorname{sg}(Sx)&=1,\label{eq:eic-sign-computation}\\
 K(0)&=1,
 &K(Sx)&=x.\label{eq:eic-countdown-computation}
\end{align}

\paragraph{Constructor laws.}
Successor is injective: applying $P$ to $Sx=Sy$ gives $x=y$ by
\eqref{eq:eic-predecessor}.  Successor never equals zero.  Indeed,
$Sx=0$ gives $1=0$ by applying $\operatorname{sg}$ and using
\eqref{eq:eic-sign-computation}; symmetry and
\eqref{eq:eic-separation} give $\bot$.  Consequently $1\ne0$, where
$A\ne B$ abbreviates $A=B\to\bot$.  None of these arguments uses
induction.

\paragraph{The sign test.}
For every $u$,
\begin{align}
 &\operatorname{sg}(u)=0\lor\operatorname{sg}(u)=1,
       \label{eq:eic-sign-range}\\
 &\operatorname{sg}(u)=0\quad\longleftrightarrow\quad u=0,
       \label{eq:eic-sign-zero}\\
 &\operatorname{sg}(u)=1\quad\longleftrightarrow\quad u\ne0.
       \label{eq:eic-sign-one}
\end{align}
To prove the range assertion, split \eqref{eq:eic-coverage}.  The branch
$u=0$ gives sign $0$, and the branch $u=S(Pu)$ gives sign $1$.
The implication $u=0\to\operatorname{sg}(u)=0$ is computation.
For its converse, the successor branch of coverage would make
$\operatorname{sg}(u)=0$ imply $1=0$, a contradiction; in that branch
explosion gives $u=0$ as well.  Similarly, if
$\operatorname{sg}(u)=1$ and $u=0$, computation gives $0=1$.
Conversely, $u\ne0$ excludes the zero branch of coverage, and the
successor branch computes the sign as $1$.

Coverage and \eqref{eq:eic-selector} also give, for all $u,v$,
\begin{equation}
 C(u,v,v)=v.
 \label{eq:eic-selector-idempotence}
\end{equation}
Both coverage branches compute to $v$, so this identity again needs no
induction.

\paragraph{The countdown test and the notation $W$.}
For every $u$,
\begin{equation}
 K(u)=0\quad\longleftrightarrow\quad u=1.
 \label{eq:eic-countdown-zero}
\end{equation}
If $u=1=S0$, \eqref{eq:eic-countdown-computation} gives $K(u)=0$.
Conversely, assume $K(u)=0$ and split coverage.  If $u=0$, then
$K(u)=1$, contradicting separation, and explosion gives $u=1$.
If $u=S(Pu)$, computation gives $K(u)=Pu$; hence $Pu=0$ and
$u=S0=1$.

If a presentation uses a description named $W$ with defining equations
\[
 W(0)=1,\qquad W(Sx)=x,
\]
then it proves $K(u)=W(u)$ for every $u$ already from coverage.
In the zero branch both sides are $1$; in the branch $u=S(Pu)$ both
sides are $Pu$.  Thus $K$ and $W$ are provably equal on the whole domain,
even if they are distinct descriptions.  This equality does not appeal to
standard-model agreement or to uniqueness of an arbitrary recurrence.

\subsection{Four elementary uses of equation induction}

Choose descriptions for addition and truncated subtraction, with recursion
on the second argument:
\begin{align}
 x+0&=x, &x+Sy&=S(x+y),\label{eq:eic-addition}\\
 d(x,0)&=x, &d(x,Sy)&=P(d(x,y)).\label{eq:eic-subtraction}
\end{align}
Using these equations and \eqref{eq:eic-induction}, we now derive
\begin{align}
 0+y&=y,\label{eq:eic-left-zero}\\
 (Sx)+y&=S(x+y),\label{eq:eic-left-successor}\\
 d(Sx,Sy)&=d(x,y),\label{eq:eic-cancel-successors}\\
 d(Sa,a)&=1.\label{eq:eic-diagonal}
\end{align}
Each induction below has a single equation as its induction formula.

For \eqref{eq:eic-left-zero}, induct on $y$.  The base is $0+0=0$.
The step is $0+Sy=S(0+y)=Sy$.

For \eqref{eq:eic-left-successor}, hold $x$ fixed and induct on $y$.
At zero both sides are $Sx$.  At a successor the induction hypothesis
gives
\[
 (Sx)+Sy=S\bigl((Sx)+y\bigr)
       =S\bigl(S(x+y)\bigr)=S(x+Sy).
\]
Combining \eqref{eq:eic-left-successor} with the defining equation for
addition gives the useful rearrangement
\begin{equation}
 (Sy)+r=y+Sr.
 \label{eq:eic-shift-successor}
\end{equation}
Indeed, both sides are $S(y+r)$.

For \eqref{eq:eic-cancel-successors}, hold $x$ fixed and induct on $y$.
The base computes as
\[
 d(Sx,S0)=P(d(Sx,0))=P(Sx)=x=d(x,0).
\]
Applying $P$ to the induction equation gives the step:
\[
 d(Sx,S(Sy))=P(d(Sx,Sy))=P(d(x,y))=d(x,Sy).
\]

Finally, induct on $a$ for \eqref{eq:eic-diagonal}.  The base is
$d(S0,0)=S0=1$.  At a successor,
\[
 d(S(Sa),Sa)=d(Sa,a)=1,
\]
using \eqref{eq:eic-cancel-successors} and the induction hypothesis.

\subsection{An equation invariant that reflects actual equality}

The remaining arithmetic ingredient is the following equation, valid for
all $x,y$:
\begin{equation}
 C\bigl(\operatorname{sg}(d(x,y)),\ x,\ y+d(x,y)\bigr)=x.
 \label{eq:eic-residual-invariant}
\end{equation}
Its meaning is that a nonzero subtraction residual reconstructs the
original argument by addition.  The selector expresses this assertion as
one literal equation, so equation induction applies directly.

Hold $x$ fixed and induct on $y$.  At $y=0$, subtraction gives residual
$x$, and \eqref{eq:eic-left-zero} gives $0+x=x$.  The desired equation
is consequently $C(\operatorname{sg}(x),x,x)=x$, which is
\eqref{eq:eic-selector-idempotence}.

For the step, assume \eqref{eq:eic-residual-invariant} at $y$ and
abbreviate
\[
 r=d(x,y),\qquad r'=Pr=d(x,Sy).
\]
These are local term abbreviations, not additional function symbols.
Split coverage for $r'$.

If $r'=0$, then $\operatorname{sg}(r')=0$, and the successor instance
of \eqref{eq:eic-residual-invariant} reduces immediately to $x=x$.

In the other branch, $r'=S(Pr')$, so constructor separation gives
$r'\ne0$.  Split coverage for $r$.  Its zero branch would imply
$r'=P0=0$, a contradiction, from which the desired step equation follows.
In its successor branch,
\[
 r=S(Pr)=Sr',\qquad \operatorname{sg}(r)=1.
\]
The induction equation therefore computes to $y+r=x$.  Meanwhile,
\eqref{eq:eic-shift-successor} gives
\[
 (Sy)+r'=y+Sr'=y+r=x.
\]
Since $r'$ is a successor, $\operatorname{sg}(r')=1$, so this is
exactly the branch selected in the successor instance of
\eqref{eq:eic-residual-invariant}.  This completes the equation induction.
All case distinctions came from coverage; no induction on an implication
or a disjunction was used.

\paragraph{Reflection at residual $1$.}
We have, for arbitrary $a,b$,
\begin{equation}
 d(Sa,b)=1\quad\longleftrightarrow\quad a=b.
 \label{eq:eic-residual-reflection}
\end{equation}
For the forward direction, substitute $x=Sa$, $y=b$ into
\eqref{eq:eic-residual-invariant}.  If $d(Sa,b)=1$, the sign is $1$
and the selector returns its second branch, giving
\[
 b+1=Sa.
\]
The defining equation for addition gives $b+1=S(b+0)=Sb$.  Thus
$Sb=Sa$, and applying $P$ gives $b=a$, hence $a=b$.
Conversely, $a=b$ and \eqref{eq:eic-diagonal} give
$d(Sa,b)=d(Sa,a)=1$.

Define the equality characteristic by the fixed description
\begin{equation}
 e(a,b)=\operatorname{sg}\bigl(K(d(Sa,b))\bigr).
 \label{eq:eic-equality-characteristic}
\end{equation}
The sign, countdown, and residual reflection lemmas give
\begin{align}
 e(a,b)=0&\quad\longleftrightarrow\quad a=b,
       \label{eq:eic-equality-reflection}\\
 e(a,b)=0&\ \lor\ e(a,b)=1.
       \label{eq:eic-equality-range}
\end{align}
In detail, $e(a,b)=0$ is equivalent to $K(d(Sa,b))=0$ by
\eqref{eq:eic-sign-zero}, which is equivalent to $d(Sa,b)=1$ by
\eqref{eq:eic-countdown-zero}, and then to $a=b$ by
\eqref{eq:eic-residual-reflection}.  The range assertion is
\eqref{eq:eic-sign-range} applied to $K(d(Sa,b))$.

The reflection property concerns the logic's actual equality throughout
the domain.  It was not assumed as part of a numerical comparison symbol,
and its proof did not restrict $a$ or $b$ to standard numerals.
It also yields decidable equality.  In the $e(a,b)=0$ branch,
\eqref{eq:eic-equality-reflection} gives $a=b$.  In the $e(a,b)=1$
branch, an assumption $a=b$ would give $e(a,b)=0$ and therefore $0=1$,
contradicting separation.  Hence
\begin{equation}
 \forall a\,\forall b\,\bigl(a=b\lor(a=b\to\bot)\bigr)
 \label{eq:eic-equality-decidable}
\end{equation}
is a theorem, not an extra hypothesis.

\subsection{Compilation of all quantifier-free formulas}

For each fixed quantifier-free formula $A(\vec p)$, construct a term
$b_A(\vec p)$ by external recursion on its finite syntax.  Take equality,
$\bot$, $\land$, $\lor$, and $\to$ as the primitive formula constructors;
negation is implication into $\bot$.  Set
\begin{align}
 b_{s=t}&=e(s,t), & b_{\bot}&=1,\label{eq:eic-boolean-atoms}\\
 b_{A\land D}&=C(b_A,b_D,1),\label{eq:eic-boolean-and}\\
 b_{A\lor D}&=C(b_A,0,b_D),\label{eq:eic-boolean-or}\\
 b_{A\to D}&=C(b_A,b_D,0).\label{eq:eic-boolean-implies}
\end{align}
Every term here is a finite composition of fixed primitive-recursive
descriptions; if desired, projections and composition give it a single
description symbol.  No internal quantification over formulas or programs
is involved.

An external induction on the syntax of $A$ proves, within the arithmetic
theory,
\begin{equation}
 \bigl(b_A=0\lor b_A=1\bigr)
 \quad\text{and}\quad
 \bigl(b_A=0\longleftrightarrow A\bigr).
 \label{eq:eic-boolean-correctness}
\end{equation}
Here are the full connective checks.  In every case $b_A=1$ implies
$\neg A$: if $A$ held, the already established implication
$A\to b_A=0$ would contradict separation.

\paragraph{Atomic formulas and falsity.}
For an equation, both claims are
\eqref{eq:eic-equality-reflection}--\eqref{eq:eic-equality-range}.
For falsity the range assertion is $1=0\lor1=1$, with its right disjunct
proved by reflexivity.  Also $1=0\to\bot$ follows from separation and
symmetry, while $\bot\to1=0$ follows by explosion.  Thus $1=0$ is
equivalent to $\bot$, as required.

\paragraph{Conjunction.}
Split $b_A=0\lor b_A=1$.  In the zero branch, $A$ holds and
$b_{A\land D}=b_D$.  Its range follows from the range of $b_D$,
and its being zero is equivalent to $D$, hence to $A\land D$ under
the known $A$.  In the one branch, $\neg A$ holds and
$b_{A\land D}=1$.  Its range is immediate.  If it were zero,
separation would give $\bot$, hence $A\land D$; conversely an
assumption $A\land D$ contradicts $\neg A$, and explosion gives
$b_{A\land D}=0$.

\paragraph{Disjunction.}
In the branch $b_A=0$, the formula $A$ holds and
$b_{A\lor D}=0$.  The range is immediate, and both $b_{A\lor D}=0$
and $A\lor D$ hold, proving their equivalence.  In the branch
$b_A=1$, we have $\neg A$ and $b_{A\lor D}=b_D$.  Its range follows
from the range of $b_D$.  If its value is zero, then $D$ and hence
$A\lor D$.  Conversely, from $A\lor D$, the $A$ alternative
contradicts $\neg A$ and yields $b_D=0$ by explosion, while the
$D$ alternative yields $b_D=0$ directly.  This proves the required
equivalence in both alternatives.

\paragraph{Implication.}
In the branch $b_A=0$, the formula $A$ holds and
$b_{A\to D}=b_D$.  Its range follows from the range of $b_D$.
If its value is zero, then $D$, which gives $A\to D$.  Conversely,
$A\to D$ and the known $A$ give $D$, hence $b_D=0$.
In the branch $b_A=1$, we have $\neg A$ and
$b_{A\to D}=0$.  Its range is immediate.  To prove $A\to D$,
assume $A$; this contradicts $\neg A$, and explosion gives $D$.
Thus both $b_{A\to D}=0$ and $A\to D$ hold in this branch, and
their equivalence follows.  This argument uses the proved two-value
range of the finite characteristic.  It does not use excluded middle
for arbitrary formulas or change the surrounding intuitionistic logic.

This completes the external syntactic induction proving
\eqref{eq:eic-boolean-correctness} for every quantifier-free formula.

\subsection{Recovery of induction and the full theory}

Let $A(n,\vec p)$ be any fixed quantifier-free formula.  Suppose an
arbitrary retained context $\Gamma$ proves its base and its induction
step:
\[
 \Gamma\vdash A(0,\vec p),
 \qquad
 \Gamma,A(z,\vec p)\vdash A(Sz,\vec p),
\]
with the usual freshness condition on $z$.  By
\eqref{eq:eic-boolean-correctness}, the base gives
$\Gamma\vdash b_A(0,\vec p)=0$.  Under the additional assumption
$b_A(z,\vec p)=0$, the same equivalence gives $A(z,\vec p)$,
the given step gives $A(Sz,\vec p)$, and the equivalence gives
$b_A(Sz,\vec p)=0$.  Apply literal-equation induction to
$b_A(n,\vec p)=0$ and then apply its reflection implication.
The result is $\Gamma\vdash A(t,\vec p)$ for every target term $t$.
Thus the theory proves the full contextual quantifier-free induction
rule, including in arbitrary quantified contexts.

\paragraph{Completion theorem.}
In the common language of fixed primitive-recursive descriptions, the
theory
\[
 \mathsf{PRE}+\mathsf{Cov}+\mathsf{Sep}+\mathsf{EI}
\]
has exactly the same sequents as ordinary first-order intuitionistic
primitive recursive arithmetic with all quantifier-free induction
instances.  The forward simulation follows from the recovered
quantifier-free induction, successor injectivity, and successor separation;
all primitive-recursive defining equations and all logical rules are
shared.  For the reverse simulation, literal-equation induction is an
instance of quantifier-free induction, and successor separation gives
$\mathsf{Sep}$.  Quantifier-free induction proves $\mathsf{Cov}$:
at zero choose its left disjunct, and at a successor use
$Sx=S(P(Sx))$, which follows from $P(Sx)=x$, for its right disjunct.

Both interpretations therefore use the full domain, defined by $x=x$,
and preserve every function symbol, equality, logical connective, and
quantifier.  Both composites have the identity isomorphism, with graph
$x=y$.  The theories are consequently identity-domain bi-interpretable
and equiconsistent.  This completion theorem is unconditional once the
literal-equation induction rule has been obtained; it has no dependency
on a classical successor-uniqueness reduction.

\section{Pairing, primitive recursion, and internal numerical composites}
\label{sec:iterator-coding}
This appendix is adapted with attribution from the progress-closure
arithmetic development~\cite{progress}. Its prerequisites have now been
proved: scalar predecessor, arithmetic reflection, and full quantifier-free
induction. Pairing is a derived proof tool; it is not part of any checker.

\subsection{Pairing from a bounded square-root search}

Once QFI is available, a familiar square-based injection gives a
particularly direct compiler. Write $a\le b$ for $d(a,b)=0$.
The usual linear-order and semiring identities follow by fixed
quantifier-free inductions from the successor, addition, and monus
equations. Multiplication itself is the scalar term
\[
                        a\cdot b=I(b,0;u.u+a).
\]
Only primitive-recursive arithmetic is involved in the following
bounded root search; there is no minimization over an unbounded domain.

For a fixed parameter $n$, put
\[
 U_n(v)=C(d((Sv)^2,n),Sv,v),\qquad
 R_n(i)=I(i,0;v.U_n(v)),\qquad q(n)=R_n(n).
\]
The update increments $v$ exactly when $(v+1)^2\le n$, and
otherwise retains it. QFI in $i$, with $n$ fixed, proves
\[
 R_n(i)\le i,\qquad R_n(i)^2\le n,\qquad
 R_n(i)=i\ \vee\ n<(S(R_n(i)))^2.              \tag{Root invariant}
\]
The base has $R_n(0)=0$. If the update increments, its test
contradicts the strict-upper-bound alternative in the third invariant.
Thus the old state equals $i$, and the new state equals $Si$;
its square bound is the increment test. If the update does not increment, the
strict upper square bound holds and the state remains fixed.
An already established strict upper square bound forces every
subsequent update to remain fixed. These are exactly the cases
needed for the three displayed invariants.

At $i=n$, either the strict upper square bound already holds or
$R_n(n)=n$, in which case $n<(n+1)^2$ supplies it. Thus
\[
                        q(n)^2\le n<(q(n)+1)^2.           \tag{Root}
\]
Monotonicity of squares shows uniqueness of a number satisfying
Root. All formulas just used are quantifier free in the scalar
language, with fixed $n$ during the induction.

Now define
\[
 \pi(a,b)=(a+b)^2+a,\qquad
 L(n)=d(n,q(n)^2),\qquad H(n)=d(q(n),L(n)).
\]
For $k=a+b$, one has $k^2\le k^2+a<(k+1)^2$ since $a\le k$.
Root uniqueness gives $q(\pi(a,b))=a+b$, and therefore
\[
                    L(\pi(a,b))=a,\qquad H(\pi(a,b))=b.
\]
Surjectivity of $\pi$ is unnecessary. The original numerical domain
is never replaced by the range of this coding; pairing is used only
to store the finite state of a fixed primitive recursion. This
standard square pairing is a convenient reconstruction here, not
a claimed new arithmetic basis.

\subsection{Compiling each primitive recursion}

Suppose the fixed target description is
\[
 f(0,\vec p)=g(\vec p),\qquad
 f(Sn,\vec p)=h(n,f(n,\vec p),\vec p).
\]
Assume by external structural induction that the earlier $g,h$
already have scalar translations $g^*,h^*$. Define
\[
 W(u,\vec p)=\pi(SL(u),h^*(L(u),H(u),\vec p)),
\]
\[
 V(n,\vec p)=I(n,\pi(0,g^*(\vec p));u.W(u,\vec p)),
 \qquad f^*(n,\vec p)=H(V(n,\vec p)).
\]
The update has only its state variable bound, with every $p_i$
fixed. Equation induction proves $L(V(n,\vec p))=n$.
The two projection equations and scalar computation then prove
both defining equations for $f^*$ internally. Initial functions
and composition translate directly. This completes the translation
of every fixed finite PR description without adding a universal
evaluator.

\subsection{Why both numerical round trips are internal}

Two fixed scalar iterators with the same seed and provably equal
updates are internally equal at all counters: use literal-equation
induction comparing their values. Their base agrees, and their
successor equations plus congruence prove the step. The same
uniqueness argument works for ordinary primitive recursion.

For an ordinary PR description, induct externally on its finite
construction. The double translation satisfies its original
defining equations, so internal equation induction identifies it
with the original operation. For a native scalar description,
earlier translated updates agree by the external induction
hypothesis; scalar uniqueness identifies the double-translated
iterator. These proofs hold at arbitrary internal counters.

Thus the domain relation $x=y$ is provably a total bijection
preserving every fixed numerical operation in both composites.
The formula constructions and round trips in
Section~\ref{sec:choice-introduction} give the corresponding formula
composites. Both interpretations preserve falsity, taking native
$0=1$ to an ordinary contradiction and ordinary $\bot$ to native
$0=1$. Finite proof translations therefore establish equiconsistency.

\section{Arithmetic strength and both ordinary interpretation composites}
\label{sec:ordinary-bridge}

We spell out the arithmetic bridge to avoid identifying standard numerical
agreement with bi-interpretability. First elaborate every fixed iterator
into an ordinary function symbol. At stage zero the symbols are $0,S$.
Given a fixed earlier-stage term $h(u,\vec p)$, introduce
$i_h(n,a,\vec p)$ with the two equations
\[
i_h(0,a,\vec p)=a,\qquad
i_h(Sn,a,\vec p)=h(i_h(n,a,\vec p),\vec p).
\]
Only finite descriptions occur; the domain contains no code for a
universal evaluator. Formal parameters are distinct when a description
is declared, while actual arguments can coincide.

Let $T_I$ be ordinary intuitionistic first-order logic with actual
equality and genuine falsity, these staged equations, contextual
induction for literal equations, and
\[
\forall x\,(x=0\vee\exists y\,x=Sy),
\qquad (0=S0)\to\bot.
\]
The second sentence is the falsity bridge. It must be stated: positive
numerical equations by themselves admit a singleton structure. Native
explosion from $0=S0$ and a translation of ordinary falsity to that
equation do not authorize silently dropping the bridge from $T_I$.

The native-to-ordinary translation is
\[
(s=t)^\circ=(s=t),\qquad
(\mathsf D_{x,y}B)^\circ=
\exists x\bigl(B^\circ(x,0)\to\forall y B^\circ(x,y)\bigr).
\]
A clause $\langle A;y.B\rangle$ becomes $A^\circ\to\forall yB^\circ$.
For the auxiliary closure used by the compact checkers,
$(\mathsf Q_yB)^\circ=B^\circ(0)\to\forall yB^\circ$.
The interpretation domain is all elements, defined by $x=x$;
equality is actual equality.

Case DI is sound because numerical coverage combines its retained base
with its successor proof. Progress DI is sound by literal-equation
induction under a temporarily assumed base. Use, DE and transport are
ordinary logical inferences under the translations. Conversely the
derived logical operations, equation induction, cases and explosion in
Section~\ref{sec:choice-introduction} recover the rules and axioms of
$T_I$. An auxiliary closure is interderivable with
$\mathsf D_{v,y}B(y)$ for unused $v$: in one direction use case DI with
the closure and its base; in the other eliminate the paired formula
and return its supplied closure. Thus formulas and context clauses have
the stated meanings in both presentations.

\subsection{The arithmetic dependency order}

The scalar update has no recursion-counter argument. We do not conceal
one in an encoding. The inherited internal proofs, reproduced in the
appendices, establish the required arithmetic in this order:
\begin{enumerate}
\item Using equation induction and numerical cases, construct a fixed
scalar predecessor term and prove $P0=0$ and $P(Sn)=n$.
Existential coverage then gives $x=0\vee x=S(Px)$: in the successor
case $x=Sy$, apply $P$ to obtain $Px=y$ and substitute.
\item Construct truncated subtraction and a numerical characteristic
of actual equality. Prove its reflection on the entire original domain.
\item Compile each fixed quantifier-free formula into its characteristic
equation and recover every contextual quantifier-free induction instance.
\item With that induction available, construct pairing and projections
by a bounded square-root search and compile every fixed general
primitive-recursive description into a scalar iterator term.
\end{enumerate}
No pairing or faithful equality detector is assumed to obtain
predecessor. No appeal to standard-model correctness replaces an
internal equality. Every induction fixes its passive parameters;
diagonal substitution occurs only into a previously proved identity.

For a target recursion
\[
f(0,\vec p)=g(\vec p),\qquad
f(Sn,\vec p)=h(n,f(n,\vec p),\vec p),
\]
let $\pi,L,H$ be the proved pairing and projections. The compiler uses
\[
V(n,\vec p)=I\bigl(n,\pi(0,g^*(\vec p));
u.\pi(SL(u),h^*(L(u),H(u),\vec p))\bigr),
\qquad f^*(n,\vec p)=H(V(n,\vec p)).
\]
Equation induction proves $L(V(n,\vec p))=n$; the two projection
identities then prove the target defining equations. The construction
is external recursion on each finite PR description, not an internal
quantification over arbitrary functions.

\begin{theorem}[Whole-domain bi-interpretability and equiconsistency]
The ordinary realization $T_I$ and full first-order iPRA are
bi-interpretable. In both theories $x=y$ defines the isomorphism from
the original domain to its double interpretation. The native Choice-DI
calculus is therefore a faithful proof presentation of a theory
bi-interpretable and equiconsistent with iPRA.
\end{theorem}
\begin{proof}
Interpret an iterator in iPRA by primitive recursion that ignores its
counter argument in the update. iPRA proves its equations, literal
equation induction, coverage and separation. In the reverse direction
use the scalar compiler. The predecessor and equality-reflection proofs
supply the successor axioms and every quantifier-free induction instance.
All ordinary logical rules are preserved.

For two iterators with the same seed and provably equal updates,
equation induction on the equality of their values proves uniqueness at
every counter, with all parameters fixed. The corresponding argument
applies to ordinary primitive recursion. Induct externally on a finite
description: its double translation satisfies its original defining
equations, so internal uniqueness identifies the two functions. This
also proves coherence when substitution changes the staged syntactic
description of an update.

The domain and equality are unchanged; hence the identity graph is
internally total, bijective and preserves every interpreted operation.
The displayed ordinary meanings of the logical macros and the paired
binder give both formula composites by external induction on finite
formula syntax. Clauses are similarly equivalent to their conditional
universal meanings. The explicit falsity bridge preserves contradiction
in both directions. Effective translations of finite proofs give
equiconsistency, without a claim that either theory proves its own
consistency.
\end{proof}

\clearpage
% Includable, self-contained research appendix. Requires amsmath, amssymb,
% amsthm and theorem/lemma/proposition/corollary environments.
% All new labels begin zcc:. Mathematical source and license are attributed below.
\section{A narrower congruence rule and its exact boundary}
\label{zcc:section}

This section strengthens the supplied report \emph{Counter Congruence
Suffices for iPRA}, specifically its \texttt{theory/counter\_theorem.tex},
in \texttt{ipra\_counter\_congruence\_2026-11-10.zip} at ProveIt commit
\texttt{76a2651601d17b34d46753f189ea026a21ce8d29}. The numerical basis,
paired logical calculus, root-replacement equality algebra, and formula
lifting are inherited. The refinement is that counter congruence needs
only a literal-zero seed and updates independent of the changing state.
A shorter representation reduces its numerical-congruence proof overhead.
A second, distinct rule permits an even shorter derivation. Finally a
rule countermodel shows why restricting the premise endpoints to
successors is not the same simplification.

These are mathematical inference results. No smaller checker implementing
these bases is claimed. The supplied source and documentation use the
MIT No Attribution License, copyright 2026 ProveIt Contributors; that
license is retained with the unchanged reference checker and producer.
Mathematical attribution is retained regardless of the license's permission.

\subsection{Complete calculus used by the refinement}

There is one numerical sort. Finite terms are generated by variables,
$0$, $S$, and $R(n,a;u,v.h)$. Only the body $h$ binds $u,v$: $u$ is the
old counter and $v$ is the previous state. Other variables are passive
parameters. Successor may instead be an abbreviation
$Sa:=R(0,a;u,v.0)$, with fresh unused $u,v$. For a conventional first-order
signature, introduce one function symbol per fixed, finite earlier-stage
body description; no variable ranges over functions or descriptions, and
there is no uniformly total evaluator. The computation schemas are
\[
 R(0,a;u,v.h)=Sa,\qquad
 R(Sn,a;u,v.h)=h[n/u,R(n,a;u,v.h)/v].
 \tag{R0, RS}\label{zcc:computation}
\]
The substitutions are simultaneous and capture-free. In particular,
the scalar-state iterator from the smallest measured checker is not being
silently used here: this $R$ body receives \emph{both} counter and state.

Object formulas and internal resources have the grammars
\[
 A::=(s=t)\mid\mathsf D_{x,y}A,
 \qquad \mathsf Q_yA(y).
\]
The variables $x,y$ of $\mathsf D$ are distinct. Contexts $\Gamma$ are
finite reusable lists of object assumptions and internal resources.
An object hypothesis has the one-node rule $\Gamma\vdash A$ when
$A\in\Gamma$. An internal lookup, written
$\Gamma\mathrel{\triangleright}\mathsf Q_yB$, also has one node when
that resource belongs to $\Gamma$; it is not an object-formula conclusion.
R0 and RS are one-node object axioms in every retained context.
The logical rules are exactly
\[
 \frac{\Gamma\mathrel{\triangleright}\mathsf Q_yB(y)
       \qquad\Gamma\vdash B(0)}
      {\Gamma\vdash B(t)}\quad(\mathsf{QE}),
\]
\[
 \frac{\Gamma\vdash\mathsf D_{x,y}B(x,y)\qquad
       \Gamma,\mathsf Q_yB(a,y)\vdash F}
      {\Gamma\vdash F}\quad(\mathsf{DE}),
\]
\[
 \frac{\Gamma,B(t,0)\vdash B(t,Sz)}
      {\Gamma\vdash\mathsf D_{x,y}B(x,y)}
      \quad(\mathsf{DI}_{0}),
 \qquad
 \frac{\Gamma,E(t,z)\vdash E(t,Sz)}
      {\Gamma\vdash\mathsf D_{x,y}E(x,y)}
      \quad(\mathsf{DI}_{1}).
\]
The progress matrix $E$ must be one literal equation, whereas $B$ may
be any object formula. In DE, $a$ is fresh for the retained context,
conclusion, and fixed parameters of $B$. In either introduction, $z$
is fresh for the context, witness $t$, and fixed matrix parameters.
Local assumptions are discharged. Weakening, exchange, contraction,
renaming, and capture-free numerical substitution are admissible
structural operations, not additional proof nodes.

The equality rule retained in every version below is root-left replacement:
\[
 \frac{\Gamma\vdash s=t\qquad\Gamma\vdash s=r}
      {\Gamma\vdash t=r}\quad(\mathsf{RL}).
 \label{zcc:root}
\]
Both displayed premises may be arbitrary proofs. Restricting its major
to a retained hypothesis is \emph{not} part of the result.
The proposed congruence rule is
\[
 \frac{\Gamma\vdash s=t}
 {\Gamma\vdash R(s,0;u,v.h(u,\vec p))=
                  R(t,0;u,v.h(u,\vec p))}
 \quad(\mathsf{ZCC}),
 \label{zcc:rule}
\]
where $v$ has no free occurrence in the displayed body, including below
nested numerical binders. Its seed is literally $0$, not merely a term
proved equal to zero. There are no primitive reflexivity, general
congruence, transport, coverage, separation, or additional induction rules.
Call the resulting calculus $\mathcal C_0$.

\subsection{Equality and cases before using ZCC}

\begin{lemma}[Root equality algebra]\label{zcc:algebra}
RL and computation derive reflexivity, symmetry and transitivity.
From proofs of $L=R$, $L=F$, and $R=G$, exactly two additional RL nodes
produce a proof of $F=G$.
\end{lemma}
\begin{proof}
Set $c_s=R(Ss,0;u,v.u)$. RS gives $c_s=s$. Using two copies as the
major and minor of RL gives $s=s$, in three nodes. An available proof
of $s=t$ and this reflexivity give $t=s$. For transitivity, reverse
$s=t$ to $t=s$ and apply RL with minor $t=r$.
For the last statement, RL applied to $L=R$ and $L=F$ gives $R=F$.
Apply RL again with that major and minor $R=G$ to obtain $F=G$.
\end{proof}

\begin{lemma}[Numerical cases]\label{zcc:cases}
For every object template $A(y)$, from $\Gamma\vdash A(0)$ and
$\Gamma\vdash A(Sz)$ one can derive $\Gamma\vdash A(t)$, where $z$
is fresh for the retained data and $t$. The construction uses neither
equality rule nor progress introduction and adds four nodes to the
sum of the two premise trees.
\end{lemma}
\begin{proof}
Choose $x$ unused in $A$. Weaken the successor proof by $A(0)$ and apply
case introduction with witness $0$, producing $\mathsf D_{x,y}A(y)$.
Eliminate it, obtaining the resource $\mathsf Q_yA(y)$ because the
witness coordinate is unused. Look up that resource and use QE with
the original base proof, weakened to this context. Close DE. The added
nodes are DI, DE, resource lookup, and QE.
\end{proof}

\subsection{Zero seed and no state dependency suffice}

Define the fixed numerical context
\[
 U(n):=R(n,0;u,v.S(Su)).
 \label{zcc:U}
\]
Its two direct computations are
\[
 U(0)=S0,\qquad U(Sz)=S(Sz).
\]
Numerical cases therefore prove
\[
                         U(n)=Sn
 \tag{U}\label{zcc:Urep}
\]
for arbitrary $n$, with arbitrary retained context. This proof has six
nodes and uses no congruence or induction rule.

\begin{theorem}[Numerical congruence from ZCC]\label{zcc:numerical}
For every fixed numerical template $F(x,\vec p)$, $\mathcal C_0$ admits
\[
 \frac{\Gamma\vdash s=t}
      {\Gamma\vdash F(s,\vec p)=F(t,\vec p)}.
\]
Its displayed construction adds exactly twenty nodes to the supplied
equality proof and introduces no progress inference.
\end{theorem}
\begin{proof}
Freshen all new binders away from the passive parameters and source terms.
Apply ZCC to $s=t$ through $U$. It is permitted because $U$ has seed zero
and its update $S(Su)$ ignores $v$. Thus $U(s)=U(t)$. The two representation
proofs \eqref{zcc:Urep} give $U(s)=Ss$ and $U(t)=St$.
Two root replacements, as in Lemma~\ref{zcc:algebra}, give
\[
                              Ss=St.
 \label{zcc:successor}
\]

For the arbitrary template $F$, set
\[
                  V_F(n):=R(Sn,0;u,v.F(u,\vec p)).
 \label{zcc:V}
\]
Its update ignores its state. A single RS instance proves
\[
                         V_F(n)=F(n,\vec p)
 \tag{V}\label{zcc:Vrep}
\]
for every $n$. No cases or equality substitution are hidden in this
computation. Apply ZCC to the derived $Ss=St$ with body $F(u,\vec p)$;
its conclusion is $V_F(s)=V_F(t)$. The two endpoint computations
\eqref{zcc:Vrep}, followed by two more root replacements, give
$F(s,\vec p)=F(t,\vec p)$.

For a source proof with $P$ nodes, the successor stage has
$(P+1)+6+6+2=P+15$ nodes: one ZCC, two six-node representations, and two
RL nodes. The final stage adds one ZCC, two one-node computations, and
two RL nodes. The total is $P+20$. Each supplied premise occurs once.
\end{proof}

The supplied report's displayed numerical-congruence construction has
$P+56$ nodes. Merely replacing its final symmetry/transitivity assembly
by the two-root-replacement lemma would improve that existing construction
to $P+44$. The present zero-seed representation gives $P+20$. Thus the
comparison separates the new representation from a shorter assembly of
old proofs. These are exact sizes of specified proof trees, not minimum
proof-size claims. Their term labels contain $F$ and may grow with its size.

\subsection{Formula transport and full consequence equivalence}

\begin{theorem}[Full transport]\label{zcc:formula}
For every native Eq/$\mathsf D$ template $A(w)$, $\mathcal C_0$ admits
\[
 \frac{\Gamma\vdash s=t\qquad\Gamma\vdash A(s)}
      {\Gamma\vdash A(t)}.
\]
Its expansion uses no additional progress induction.
\end{theorem}
\begin{proof}
Use external induction on the finite number of paired binders in $A$,
uniformly in all retained contexts. For $A(w)\equiv P(w)=Q(w)$, obtain
$P(s)=P(t)$ and $Q(s)=Q(t)$ by Theorem~\ref{zcc:numerical}, then use
root equality algebra to join them with $P(s)=Q(s)$.

For $A(w)=\mathsf D_{x,y}B(w,x,y)$, eliminate $A(s)$ at a fresh $a$.
The retained resource is $\mathsf Q_yB(s,a,y)$. Introduce the desired
$A(t)$ by case introduction with witness $a$ and fresh $z$. Its local
assumption is $B(t,a,0)$. Symmetry gives $t=s$, so the external induction
hypothesis transports that assumption backwards to $B(s,a,0)$.
QE then supplies $B(s,a,Sz)$. A forward transport by the same external
induction hypothesis gives $B(t,a,Sz)$. Close case introduction and DE.
Both recursive calls concern the strictly smaller matrix $B$. Freshness
conditions are met by renaming. This is structural reasoning about
finite formulas, not compound arithmetic induction.
\end{proof}

\begin{corollary}[Identity-language equivalence]\label{zcc:equivalence}
Replacing the supplied unrestricted counter-congruence rule by ZCC leaves
every native sequent derivable exactly when it was derivable before.
The same holds when comparing with the supplied unrestricted
object-transport calculus. Therefore the inherited whole-domain
bi-interpretation and equiconsistency with full first-order iPRA transfer.
\end{corollary}
\begin{proof}
Every ZCC inference is an old counter-congruence inference. Conversely,
Theorem~\ref{zcc:numerical}, with context $R(x,a;u,v.h)$, recovers every
old counter-congruence instance, including those with a state-dependent
body and arbitrary seed. Theorem~\ref{zcc:formula} recovers full object
transport. In the other direction, RL is full transport at the left
root, while ZCC follows by transporting one occurrence in a reflexivity
equation. The transformations recursively preserve end sequents and
freshness. The arithmetic and logical iPRA bridge is the attributed
supplied bridge; this refinement introduces no new domain or equality.
\end{proof}

\subsection{A different rule with only five nodes of overhead}

Define a separate primitive rule
\[
 \frac{\Gamma\vdash s=t}
 {\Gamma\vdash R(Ss,0;u,v.h(u,\vec p))=
                  R(St,0;u,v.h(u,\vec p))}
 \quad(\mathsf{ZCC}^{+}),
 \label{zcc:wrapped}
\]
again with no free occurrence of $v$ in $h$. Its premise endpoints are
arbitrary terms; the rule itself wraps them in successors. This is a
changed rule, \emph{not} ZCC with a side condition on its premise.

\begin{proposition}\label{zcc:wrapped-theorem}
Replacing ZCC by $\mathsf{ZCC}^{+}$ also gives the full native consequence
relation. Arbitrary numerical congruence now has a $P+5$-node construction
from a $P$-node source equality, without cases or progress induction.
\end{proposition}
\begin{proof}
For arbitrary $F$, apply $\mathsf{ZCC}^{+}$ using body $F(u,\vec p)$.
This directly gives $V_F(s)=V_F(t)$. Two RS endpoint computations and
two RL nodes give $F(s)=F(t)$, adding $1+1+1+2=5$ nodes.
Formula lifting is unchanged. Conversely, Theorem~\ref{zcc:numerical}
in $\mathcal C_0$ derives $\mathsf{ZCC}^{+}$ by choosing context
$R(Sx,0;u,v.h)$. Thus the two full calculi are equivalent.
\end{proof}

This factors numerical congruence through a fixed successor-computation
shape. It is not claimed to remove the cost of numerical binding or
substitution: that cost remains in the computation rules. No source-size
figure is assigned to this alternative calculus.

\subsection{Why successor-only premise endpoints are insufficient}

Consider instead retaining ZCC only when its \emph{premise} is syntactically
$Sp=Sq$, so its conclusion is $R(Sp,0;h)=R(Sq,0;h)$. Keep all other rules
unchanged and supply no rule wrapping arbitrary premise endpoints.

\begin{proposition}[A rule countermodel]\label{zcc:countermodel}
This successor-premise-only restriction does not derive full numerical
congruence and does not preserve the old native consequence relation.
In particular it does not derive
\[
                         0=S0\ \vdash\ S0=S(S0).
\]
\end{proposition}
\begin{proof}
Use the standard numerical domain $\mathbb N$. Interpret $0,S,R$ by
the ordinary zero, successor, and successor-seeded primitive recursion.
Interpret a native equation $s=t$ by
\[
 e(s^{\mathbb N},t^{\mathbb N}),
 \qquad e(m,n)\iff m=n\ \text{or}\ m,n\in\{0,1\}.
\]
For avoidance of ambiguity, this is a \emph{rule interpretation} of the
native equation atom; it is not asserted to be a model of ordinary
first-order logic with that atom as actual equality. Interpret resources
and paired formulas by their ordinary classical meanings
\[
 (\mathsf Q_yB)^\circ=B^\circ(0)\to\forall yB^\circ(y),\qquad
 (\mathsf D_{x,y}B)^\circ=
 \exists x\bigl(B^\circ(x,0)\to\forall yB^\circ(x,y)\bigr).
\]
Classical truth suffices for a rule-soundness countermodel of derivability.

The relation $e$ is an equivalence relation, so RL is sound. Every
computation equation has actually equal values and therefore is $e$-true.
For the restricted congruence rule, both counter values are positive.
On positive numbers $e$ agrees with actual equality. Consequently its
true premise implies identical actual counters, and the recursor
values agree. In fact this argument would allow arbitrary seeds and
updates whenever both premise counters are successors.

Hypothesis use, resource lookup, QE and DE are sound by their displayed
meanings. Case introduction is sound by standard numerical coverage.
Progress introduction is sound because ordinary induction on standard
$\mathbb N$ holds for the interpreted predicate, including predicates
whose equation atoms use $e$; only its actual truth set is relevant.
Freshness and substitution have their usual numerical meanings.
Thus every rule of the restricted calculus preserves truth in this
interpretation.

The premise $0=S0$ is true since $e(0,1)$ holds, while the conclusion is
false since $e(1,2)$ does not. Hence that sequent is underivable. The old
calculus derives it by successor congruence, proving the claimed strict
loss of native consequences.
\end{proof}

This countermodel rules out the proposed identity-language restriction.
It is not a theorem excluding every possible alternate interpretation,
changed equality convention, or enlarged axiom system.

\subsection{Evidence, overlap, and limits}

The source screen recursively inspected 32 distinct ZIP payloads from the
17 current archives, with 867 distinct text files and 443 parsed Python
ASTs. It found the expected old state-free-context remark but no explicit
zero-seed congruence restriction, wrapped rule, or successor-premise
countermodel. Pattern screening does not decide semantic duplication or
establish literature-wide priority. The result is attributed as a
refinement of the existing counter-congruence argument, not an invention
of its arithmetic or logical basis.

The supporting script \texttt{validate\_zero\_seed.py} preserves the
unchanged upstream readable checker and intercepts every CC call to
require literal-zero seed and absence of the outer state slot, correctly
tracking nested numerical binders. It constructs three symbolic proofs
(successor, initializer, and nested-update congruence), each with a one-node
source and exactly 21 nodes. A separate, explicitly coded wrapped-rule
branch verifies the corresponding three six-node proofs. All six pass.
Controls with a nonzero seed, a direct state occurrence, and a state
occurrence below another recursor are rejected. A finite check of the
countermodel relation is supplied only as a sanity check; the universal
rule-soundness argument is the preceding proof.

These artifacts verify concrete proof trees and side conditions. They are
not Lean/Rocq formalizations, universal correctness proofs of a new checker,
or measurements below the 558-byte current frontier. The numerical proof
bounds do not bound the size of formula-transport lifting, serialized term
labels, or runtime. No absolute simplest axiomatization or significant
checker-size breakthrough follows from this appendix alone.


\begin{thebibliography}{9}
\bibitem{snapshot} Vladimir Reshetnikov, \emph{ProveIt, current draft corpus},
commit \file{76a2651601d17b34d46753f189ea026a21ce8d29},
10 November 2026.\newline
\url{https://github.com/VladimirReshetnikov/ProveIt/tree/76a2651601d17b34d46753f189ea026a21ce8d29/docs/drafts}.
\bibitem{certified} \emph{A Certified Zero Removes a Computation Opcode},
\file{ipra_certified_zero_2026-11-10.zip}, in~\cite{snapshot}.
Source of the selected comparison programs, certified-zero computation
repair, named emitter, and reproduced complete arithmetic/logical bridge.
\bibitem{arity} \emph{An Arity Calculus for iPRA},
\file{ipra_arity_calculus_2026-11-10.zip}, in~\cite{snapshot}.
Source of the arity encoding and sorted-opening argument, as inherited
by~\cite{certified}.
\bibitem{forward} Supplied actual-equality reflection proof,
\file{reference/equality_reflection.tex} in~\cite{certified}, with its
attribution to the earlier forward-left development.
\bibitem{erasure} Supplied scalar predecessor proof,
\file{reference/scalar_predecessor.tex} in~\cite{certified}, with its
attribution to the earlier scalar and unparenthesized one-hole developments.
\bibitem{progress} Supplied pairing and primitive-recursion proof,
\file{reference/primitive_recursion.tex} in~\cite{certified}, with its
attribution to the earlier progress-closure development.
\bibitem{counter} \emph{Counter Congruence Suffices for iPRA},
\file{ipra_counter_congruence_2026-11-10.zip}, in~\cite{snapshot}.
Source of the different $R$ calculus refined in Appendix~\ref{zcc:section}.
\end{thebibliography}
\end{document}
